\exercisesection{}

\begin{enumerate}
\def\labelenumi{\arabic{enumi})}
\item
  设 G 为局部紧群。设 \(\pi\) 为 G 在复希尔伯特空间 E 中的酉表示，且 \(x \in \mathrm{E}\) 。我们有 \(\pi(g)x = x\) （对所有 \(g \in \mathrm{G}\) ），当且仅当正定函数 \(g \mapsto \langle x | \pi(g)x \rangle\) 为常数。
\item
  设 G 为局部紧群。对 G 上的每个正定函数 \(\varphi\) ，我们用 \((\pi_{\varphi}, x_{\varphi})\) 表示 \(\varphi\) 的一个循环希尔伯特表示。
\end{enumerate}

\begin{enumerate}
\def\labelenumi{\alph{enumi})}
\item
  设 \(\varphi_1, \varphi_2\) 与 \(\varphi_3\) 为 G 上的正定函数。若 \(\varphi_3 = \varphi_1 + \varphi_2\) ，则存在从 \(\pi_{\varphi_1}\) 到 \(\pi_{\varphi_3}\) 的单射 G-态射。
\item
  若 \(\varphi\) 是非零的常值正定函数，则 \(\pi_{\varphi}\) 是一维的平凡表示。
\item
  对 G 在希尔伯特空间 E 中的每个酉表示 \(\pi\) 以及每个 \(x \in \mathrm{E}\) ，存在从 \(\pi_{\varphi}\) 到 \(\pi\) 的单射 G-态射，其中 \(\varphi(g) = \langle x \mid \pi(g)x \rangle\) 。
\item
  设 H 为 G 的闭子群，\(\varphi\) 为 G 上的正定函数。用 \(\psi\) 表示 \(\varphi\) 在 H 上的限制；它是 H 上的正定函数，且存在从 \(\pi_{\psi}\) 到 \(\operatorname{Res}_{\mathrm{H}}^{\mathrm{H}}(\pi)\) 的单射 H-态射。
\end{enumerate}

\begin{enumerate}
\def\labelenumi{\arabic{enumi})}
\setcounter{enumi}{2}
\tightlist
\item
  \begin{enumerate}
  \def\labelenumii{\alph{enumii})}
  \tightlist
  \item
  给出一个对合巴拿赫代数的例子，它具有一个不连续的正线性泛函。
  \end{enumerate}
\end{enumerate}

\begin{enumerate}
\def\labelenumi{\alph{enumi})}
\setcounter{enumi}{1}
\tightlist
\item
  给出一个对合巴拿赫代数 A、A 上一个连续的正线性泛函 \(\lambda\) ，以及希尔伯特实现 \((\mathrm{E}_{1}, e_{1}, \varphi_{1})\) 与 \((\mathrm{E}_{2}, e_{2}, \varphi_{2})\) （均为 \(\lambda\) 的实现）的例子，使得 \((\mathrm{E}_{1}, e_{1}, \varphi_{1})\) 是循环希尔伯特实现，并且 V, p.~440 命题 6 中唯一的线性映射 u 不满足 \(u(e_{1}) = e_{2}\) 。
\end{enumerate}

\begin{enumerate}
\def\labelenumi{\arabic{enumi})}
\setcounter{enumi}{3}
\tightlist
\item
  设 A 与 B 为星代数。从 A 到 B 的线性映射 u 称为正的，如果 \(u(\mathrm{A}_{+}) \subset \mathrm{B}_{+}\) ；称为完全正的，如果对每个整数 \(k \in N\) ，映射 \(u_{k} = 1_{C^{k}} \otimes u: C^{k} \otimes A \to C^{k} \otimes B\) 是正的。
\end{enumerate}

\begin{enumerate}
\def\labelenumi{\alph{enumi})}
\tightlist
\item
  设 A 为含幺星代数，E 与 F 为复希尔伯特空间，v 为从 E 到 F 的线性映射，\(\varphi\) 为从 A 到 \(\mathcal{L}(\mathrm{F})\) 的对合代数的保幺态射，使得 \(\|\varphi\|^{2} = \|\varphi(1)\|\) 。映射 \(u: A \to \mathcal{L}(E)\) （由 \(u(a) = v^{*} \circ \varphi(a) \circ v\) 定义）是完全正的。
\end{enumerate}

本习题的目标是证明这一断言的逆命题（“斯廷斯普龙定理”）。设 A 为含幺星代数，E 为复希尔伯特空间，\(u \colon \mathrm{A} \to \mathcal{L}(\mathrm{E})\) 为完全正线性映射。

\begin{enumerate}
\def\labelenumi{\alph{enumi})}
\setcounter{enumi}{1}
\tightlist
\item
  设 \(F = A \otimes E\) 。F 上存在一个正的半双线性型 q，使得
\end{enumerate}

\[
q (a \otimes x, b \otimes y) = \langle x | \varphi (a ^ {*} b) y \rangle
\]

对所有 \((a,b)\in \mathrm{A}^2\) 与 \((x,y)\in \mathrm{E}^2\) 成立。

设 \(\widehat{\mathbf{F}}\) 是 \(\mathbf{F}\) 关于半双线性型 \(q\) 的分离完备希尔伯特空间。

\begin{enumerate}
\def\labelenumi{\alph{enumi})}
\setcounter{enumi}{2}
\tightlist
\item
  设 \(b \in A\) 。F 上使 \(a \otimes x \mapsto ba \otimes x\) （对所有 \(a \in A\) 与 \(x \in E\) ）成立的自同态，通过过渡到商诱导出 \(\varphi(b)\) （\(\widehat{F}\) 的一个自同态），满足 \(\|\varphi(b)\|\leqslant\|b\|\) ；映射 \(\varphi\colon A\to\mathcal{L}(F)\) 是对合代数的保幺态射。
\end{enumerate}

\begin{enumerate}
\def\labelenumi{\alph{enumi})}
\setcounter{enumi}{3}
\tightlist
\item
  设 \(v\) 为从 E 到 \(\widehat{\mathrm{F}}\) 的线性映射，使得 \(v(x)\) 是 \(1 \otimes x\) 所在的类（对所有 \(x \in \mathrm{E}\) ）。我们有 \(\|v\|^2 = \| \varphi(1)\|\) ，且 \(u(a) = v^* \circ \varphi(a) \circ v\) 对所有 \(a \in \mathrm{A}\) 成立。
\end{enumerate}

\begin{enumerate}
\def\labelenumi{\arabic{enumi})}
\setcounter{enumi}{4}
\tightlist
\item
  设 X 为分离拓扑空间。连续函数 \(f: X \times X \to R\) 称为条件负定型的，如果它满足下列条件
\end{enumerate}

\begin{enumerate}
\def\labelenumi{(\roman{enumi})}
\item
  对每个 \(x \in X\) ，有 \(f(x, x) = 0\) ；
\item
  对所有 \(x\) 与 \(y\) （属于 \(X\) ），我们有 \(f(x,y) = f(y,x)\) ；
\item
  对每个整数 \(n \in N\) 、X 中元素的每个族 \((x_{i})_{0 \leqslant i \leqslant n}\) 以及实数的每个族 \((t_{i})_{0 \leqslant i \leqslant n}\) （满足
\end{enumerate}

\[
\sum_ {i = 1} ^ {n} t _ {i} = 0,
\]

），我们有

\[
\sum_ {i = 0} ^ {n} \sum_ {j = 0} ^ {n} t _ {i} t _ {j} f (x _ {i}, x _ {j}) \leqslant 0.
\]

我们用 Noy \(_{-}\) (X) 表示 X × X 上条件负定型函数构成的集合。

\begin{enumerate}
\def\labelenumi{\alph{enumi})}
\item
  设 E 为实希尔伯特空间，\(\varphi\colon\mathrm{X}\to\mathrm{E}\) 为连续映射。函数 \(f(x,y)=\|\varphi(x)-\varphi(y)\|^{2}\) 是条件负定型的。
\item
  设 \(f\in\mathrm{Noy}_{-}(\mathrm{X})\) 且 \(x_{0}\in\mathrm{X}\) 。存在一个实希尔伯特空间 E 与一个连续映射 \(\varphi\colon\mathrm{X}\to\mathrm{E}\) ，使得
  \begin{enumerate}
  \def\labelenumii{\alph{enumii})}
  \tightlist
  \item
    对每个 \((x,y)\in \mathrm{X}\times \mathrm{X}\) ，有 \(f(x,y) = \| \varphi (x) - \varphi (y)\| ^2\) ；
  \item
    族 \((\varphi(x)-\varphi(x_{0}))_{x\in\mathrm{X}}\) 在 E 中是完全的。
  \end{enumerate}
\end{enumerate}

我们称 \((\mathrm{E},\varphi)\) 为 f 的一个循环希尔伯特表示。

\begin{enumerate}
\def\labelenumi{\alph{enumi})}
\setcounter{enumi}{2}
\item
  设 \(f \in \mathrm{Noy}_{-}(\mathrm{X})\) 且 \(x_0 \in \mathrm{X}\) 。设 \((\mathrm{E}, \varphi)\) 与 \((\mathrm{E}', \varphi')\) 为 \(f\) 的循环希尔伯特表示。存在唯一的仿射等距 \(u: \mathrm{E} \to \mathrm{E}'\) ，使得 \(\varphi' = u \circ \varphi\) 。
\item
  集合 \(\mathrm{Noy}_{-}(\mathrm{X})\) 是 \(\mathcal{C}(\mathrm{X}\times \mathrm{X};\mathbf{R})\) 中的一个凸锥。对每个普遍正核 \(g\in \mathrm{Noy}_{+}(\mathrm{X})\) （\(g\) 取实值），映射 \(f(x,y) = g(x,x) - g(x,y)\) 属于 \(\mathrm{Noy}_{-}(\mathrm{X})\) 。
\item
  设 \(\mathfrak{F}\) 为 \(\mathrm{Noy}_{-}(\mathrm{X})\) 上逐点收敛于 \(f\colon \mathrm{X}\times \mathrm{X}\to \mathbf{R}\) 的滤子。若 \(f\) 连续，则 \(f\in \mathrm{Noy}_{-}(\mathrm{X})\) 。
\item
  设 \(f: X \times X \to R\) 为满足上述条件 (i) 与 (ii) 的连续函数。设 \(x_{0} \in X\)。设 \(g: X \times X \to R\) 满足 \(g(x, y) = f(x, x_{0}) + f(y, x_{0}) - f(x, y)\) 。则 \(f \in \text{Noy}_{-}(X)\) 当且仅当 \(g \in \text{Noy}_{+}(X)\) 。
\item
  假设 X 非空。设 \(f: \mathrm{X} \times \mathrm{X} \to \mathbf{R}\) 为满足上述条件 (i) 与 (ii) 的连续函数。则 \(f \in \mathrm{Noy}_{-}(\mathrm{X})\) 当且仅当 \(\exp(-tf) \in \mathrm{Noy}_{+}(\mathrm{X})\) 对所有 \(t \in \mathbf{R}_{+}\) 成立（“舍恩贝格定理”）。（为证 \(\exp(-f)\) 在 \(f \in \mathrm{Noy}_{-}(\mathrm{X})\) 时是正定型的，固定 \(x_0 \in \mathrm{X}\) ；证明
\end{enumerate}

\[
g (x, y) = \exp (f (x, x _ {0}) + f (y, x _ {0}) - f (x, y))
\]

以及

\[
h (x, y) = \exp (- f (x, x _ {0})) \exp (- f (y, x _ {0}))
\]

都属于 Noy \(_{+}\) (X)。）

若 G 是拓扑群，我们称函数 \(\varphi \in \mathcal{C}(\mathrm{G})\) 是负定型的，如果函数 \(f\) 是 \(\mathrm{G} \times \mathrm{G}\) 上由 \(f(g,h) = \varphi(h^{-1}g)\) 定义的条件负定核。

\begin{enumerate}
\def\labelenumi{\arabic{enumi})}
\setcounter{enumi}{5}
\tightlist
\item
  设 G 为离散拓扑群。设 A 为 \(\mathcal{L}(\ell^{2}(\mathrm{G}))\) 中由左正则表示 \(\lambda\) 的像生成的对合子代数。这是一个含幺星代数。
\end{enumerate}

对每个正定型函数 \(\varphi\in\mathrm{Pos}(\mathrm{G})\) ，存在唯一的从 A 到 A 的完全正线性映射 \(u_{\varphi}\) （参见习题 4），使得 \(u_{\varphi}(\boldsymbol{\lambda}(g))=\varphi(g)\boldsymbol{\lambda}(g)\) 对所有 \(g\in\mathrm{G}\) 成立。（利用 \(\varphi\) 的一个循环希尔伯特表示。）

\begin{enumerate}
\def\labelenumi{\arabic{enumi})}
\setcounter{enumi}{6}
\tightlist
\item
  设 I 为有限集，\(G = F(I)\) 为在 I 上构造的自由群（A, I, p.~84）。设 A 为 \(\mathcal{L}(\ell^{2}(G))\) 中由左正则表示 \(\lambda\) 的像生成的对合子代数。
\end{enumerate}

\begin{enumerate}
\def\labelenumi{\alph{enumi})}
\item
  对每个 \(g \in G\) ，用 \(|g|\) 表示 g 的唯一既约分解的长度（A, I, p.~146, 习题 19）。映射 \(d(g, h) = |h^{-1}g|\) 是 G 上的一个距离。
\item
  设 \(\lambda \in \mathbf{R}_{+}^{*}\) 。G 上由 \(\psi_{\lambda}(g) = \exp(-\lambda|g|)\) 定义的函数是正定函数。（证明函数 \(g \mapsto |g|\) 是负定型函数，参见习题~5。）
\item
  对每个整数 \(k \in N\) ，设 \(G_{k}\) 为由 \(g \in G\) （满足 \(|g| = k\) ）组成的集合，并设 \(\varphi_{k}\) 为 \(G_{k}\) 的特征函数。设 f 与 g 为 G 上的函数，其支集分别含于 \(G_{k}\) 与 \(G_{l}\) 。对每个 \(m \in N\) ，我们有 \((f * g)\varphi_{m} = 0\) ，除非 \(k + l - m\) 为偶数且 \(|k - l| \leqslant m \leqslant k + l\) ；在这一情形，我们有
\end{enumerate}

\[
\left\| (f * g) \varphi_ {m} \right\| \leqslant \| f \| \| g \|.
\]

\begin{enumerate}
\def\labelenumi{\alph{enumi})}
\setcounter{enumi}{3}
\tightlist
\item
  设 f 为 G 上具有有限支集的函数。我们有
\end{enumerate}

\[
\| \boldsymbol {\lambda} (f) \| \leqslant \sum_ {n \in \mathbf {N}} (n + 1) \| f \varphi_ {n} \|,
\]

以及

\[
\left\| \boldsymbol {\lambda} (f) \right\| \leqslant 2 \left(\sum_ {g \in G} (1 + | g |) ^ {4} | f (g) | ^ {2}\right) ^ {1 / 2}.
\]

\begin{enumerate}
\def\labelenumi{\alph{enumi})}
\setcounter{enumi}{4}
\item
  设 \(\psi\) 为 G 上的函数，使得 \(M = \sup_{g \in G} |\psi(g)| (1 + |g|)^2\) 有限。存在唯一的从 A 到 A 的线性映射 \(u_\psi\) ，使得 \(u_\psi(\boldsymbol{\lambda}(g)) = \psi(g)\boldsymbol{\lambda}(g)\) 对所有 \(g \in G\) 成立。我们有 \(\|u_\psi\| \leqslant 2M\) 。
\item
  设 X 为 G 上由带有限支集的函数 \(\psi\) 组成的集合，其中 \(\|u_{\psi}\|\leqslant1\) 。存在 X 上的一个滤子 \(\mathfrak{F}\) ，使得
\end{enumerate}

\[
\lim _ {\psi , \mathfrak {F}} u _ {\psi} = 1 _ {\mathrm{A}}
\]

在赋予逐点收敛拓扑的空间 \(\mathcal{L}(\mathrm{A})\) 中。特别地，A 的底巴拿赫空间具有逼近性质。（考虑函数 \(\psi_{\lambda,n}(g) = (\varphi_0 + \dots + \varphi_n)\psi_\lambda\) ，其中 \(\lambda \in \mathbf{R}_+^*\) ，\(n \in \mathbf{N}\) ，并利用上一习题。）

\begin{enumerate}
\def\labelenumi{\arabic{enumi})}
\setcounter{enumi}{7}
\tightlist
\item
  我们把群 \(\mathbf{R}\) 与它的对偶群通过映射 \((x,y)\mapsto\) \(\exp (i\pi xy)\) 等同起来。设 \(\alpha \in \mathbf{R}_{+}\) 。
\end{enumerate}

\begin{enumerate}
\def\labelenumi{\alph{enumi})}
\item
  \(\mathbf{R}\) 上存在其傅里叶变换为函数 \(f_{\alpha}(x) = \exp(-|x|^{\alpha})\) 的正测度，当且仅当 \(\alpha \in ]0,2]\) 。在这一情形，我们把这一测度记作 \(\mu_{\alpha}\) ，它是唯一的；称之为参数 \(\alpha\) 的稳定测度。
\item
  测度 \(\mu_{\alpha}\) 的总质量为 1；R 上的恒等函数 f 关于 \(\mu_{\alpha}\) 可积，当且仅当 \(\alpha = 2\) 。
\end{enumerate}

\begin{enumerate}
\def\labelenumi{\arabic{enumi})}
\setcounter{enumi}{8}
\tightlist
\item
  本习题的目标是给出 V, p.~455 定理 5 的另一个证明。设 G 为局部紧交换拓扑群，带有哈尔测度 \(\mu\) 。我们用 \(\widehat{f}\) 表示 \(f \in \mathrm{L}^1(\mathrm{G})\) 的傅里叶变换，并用 \(\mathcal{F}(\mathrm{G}) \subset \mathcal{C}_0(\widehat{\mathrm{G}})\) 表示 \(\mathrm{L}^1(\mathrm{G})\) 上傅里叶变换的像。
\end{enumerate}

\begin{enumerate}
\def\labelenumi{\alph{enumi})}
\item
  空间 \(\mathcal{F}(\mathrm{G})\) 是 \(\mathcal{C}_0(\mathrm{G})\) 的稠密子空间。
\item
  对每个 \(f \in \mathrm{L}^{1}(\mathrm{G})\) ，f 在巴拿赫代数 \(\mathrm{L}^{1}(\mathrm{G})\) 中的谱半径等于 \(\|\widehat{f}\|_{\infty}\) 。
\item
  设 \(\varphi\in\mathrm{Pos}(\mathrm{G})\) 。它是 \(\mathrm{G}\) 上一致连续且有界的函数；对每个 \(f\in\mathrm{L}^{1}(\mathrm{G})\) ，我们有
\end{enumerate}

\[
\int_ {G} \int_ {G} f (x) \overline {{f (y)}} \varphi (x - y) d (\mu \otimes \mu) (x, y) \geqslant 0.
\]

\begin{enumerate}
\def\labelenumi{\alph{enumi})}
\setcounter{enumi}{3}
\tightlist
\item
  设 \(\lambda\) 为 \(\mathrm{L}^{1}(\mathrm{G})\) 上由 \(\varphi\) 确定的连续线性型，并令 \(b(f,g)=\lambda(f^{*}*g)\) （对所有 \((f,g)\in\mathrm{L}^{1}(\mathrm{G})\times\mathrm{L}^{1}(\mathrm{G})\) ）。映射 \(b\) 是 \(\mathrm{L}^{1}(\mathrm{G})\) 上的一个正埃尔米特型。
\end{enumerate}

\begin{enumerate}
\def\labelenumi{\alph{enumi})}
\setcounter{enumi}{4}
\tightlist
\item
  对每个 \(f \in \mathrm{L}^1(\mathrm{G})\) ，我们有 \(|\lambda(f)|^2 \leqslant b(f, f)\) 。
\end{enumerate}

\begin{enumerate}
\def\labelenumi{\alph{enumi})}
\setcounter{enumi}{5}
\tightlist
\item
  设 \(f \in \mathrm{L}^{1}(\mathrm{G})\) ，\(g = f^{*} * f \in \mathcal{C}_{b}(\mathrm{G})\) 。定义 \(g_{1} = g\) 与 \(g_{n+1} = g_{n} * g\) （对每个整数 \(n \geqslant 1\) ）。对每个整数 \(n \geqslant 1\) ，我们有 \(|\lambda(f)| \leqslant \|g_{n}\|_{1}^{2^{-n}}\) 。
\end{enumerate}

\begin{enumerate}
\def\labelenumi{\alph{enumi})}
\setcounter{enumi}{6}
\tightlist
\item
  对每个 \(f \in \mathrm{L}^{1}(\mathrm{G})\) ，我们有 \(|\lambda(f)| \leqslant \| \widehat{f} \|_{\infty}\) 。存在连续线性映射 \(\widehat{\lambda}\) （定义在 \(\mathcal{F}(\mathrm{G})\) 上），使得 \(\lambda(f) = \widehat{\lambda}(\widehat{f})\) 对所有 \(f \in \mathrm{L}^{1}(\mathrm{G})\) 成立。
\end{enumerate}

\begin{enumerate}
\def\labelenumi{\alph{enumi})}
\setcounter{enumi}{7}
\tightlist
\item
  存在有界测度 \(\nu\) （在 \(\widehat{G}\) 上），使得 \(\lambda(f) = \nu(\widehat{f})\) 对每个 \(f \in \mathrm{L}^{1}(\mathrm{G})\) 成立；测度 \(\nu\) 是正的，并且有
\end{enumerate}

\[
\varphi (x) = \int_ {\widehat {\mathrm{G}}} \langle \chi , x \rangle d \nu (\chi)
\]

对所有 \(x \in G\) 成立。

\begin{enumerate}
\def\labelenumi{\arabic{enumi})}
\setcounter{enumi}{9}
\tightlist
\item
  \begin{enumerate}
  \def\labelenumii{\alph{enumii})}
  \tightlist
  \item
    设 \(\mathrm{G} = \mathbf{R}\) 。集合 \(\operatorname{Pos}_1(\mathrm{G})\) 在弱拓扑下不是紧的。
  \end{enumerate}
\end{enumerate}

\begin{enumerate}
\def\labelenumi{\alph{enumi})}
\setcounter{enumi}{1}
\tightlist
\item
  设 G = R/Z。\(\operatorname{Pos}_{\leqslant1}(\mathbf{G})\) 上的弱拓扑与紧收敛拓扑不一致。
\end{enumerate}

\begin{enumerate}
\def\labelenumi{\arabic{enumi})}
\setcounter{enumi}{10}
\tightlist
\item
  设 G 为局部紧群。
\end{enumerate}

设 \(\varrho\) 与 \(\pi\) 为 G 的酉表示。我们称 \(\varrho\) 弱包含于 \(\pi\) ，并记作 \(\varrho \leqslant \pi\) ，如果 \(\varrho\) 的对角矩阵系数所成的集合包含于 \(\pi\) 的对角矩阵系数的有限和所成的集合在赋予紧收敛拓扑的空间 \(\mathcal{C}(\mathrm{G})\) 中的闭包内。

\begin{enumerate}
\def\labelenumi{\alph{enumi})}
\item
  关系“\(\varrho\) 与 \(\pi\) 是 G 的酉表示且 \(\varrho \leqslant \pi\)”是一个预序关系。若 \(\varrho'\) （相应地 \(\pi'\) ）同构于 \(\varrho\) （相应地 \(\pi\) ），则 \(\varrho \leqslant \pi\) 当且仅当 \(\varrho' \leqslant \pi'\) 。
\item
  若 \(\varrho\) 同构于 \(\pi\) 的一个子表示，则 \(\varrho \leqslant \pi\) 。若 I 与 J 是非空有限集，且若 \(\varrho \leqslant \pi\) ，则 \(\bigoplus_{i \in I} \varrho \leqslant \bigoplus_{j \in J} \pi\) 。
\item
  设 X 为 \(\varrho\) 的空间 E 中生成 \(\pi\) 的一个子集。我们有 \(\varrho \leqslant \pi\) ，当且仅当每个矩阵系数 \(g \mapsto \langle x | \varrho(g)x \rangle\) （带 \(x \in X\) ）都属于 \(\pi\) 的对角矩阵系数的有限和所成的集合在赋予紧收敛拓扑的空间 \(\mathcal{C}(\mathrm{G})\) 中的闭包。（证明 E 中满足这一性质的元素所成的集合 Y 是 E 的一个闭向量子空间。）
\item
  设 H 为 G 的闭子群。设 \(\pi\) 与 \(\varrho\) 为 G 的酉表示，使得 \(\varrho \leqslant \pi\) ；则我们有 \(\operatorname{Res}_{\mathrm{H}}^{\mathrm{G}}(\varrho) \leqslant \operatorname{Res}_{\mathrm{H}}^{\mathrm{G}}(\pi)\) 。
\item
  设 H 为 G 的闭子群，\(\nu\) 为 H 上的哈尔测度。设 \(\pi\) 与 \(\varrho\) 为 H 的酉表示，使得 \(\varrho \leqslant \pi\) ；则我们有 \(\operatorname{Ind}_{\mathrm{H}}^{\mathrm{G}}(\varrho) \leqslant \operatorname{Ind}_{\mathrm{H}}^{\mathrm{G}}(\pi)\) 。（令 \(\varrho' = \operatorname{Ind}_{\mathrm{H}}^{\mathrm{G}}(\varrho)\) ，\(\pi' = \operatorname{Ind}_{\mathrm{H}}^{\mathrm{G}}(\pi)\) ；利用 (c)，化归为逼近形如 \(g \mapsto \langle f | \varrho'(g)f \rangle\) 的矩阵系数，其中 f 形如
\end{enumerate}

\[
f (x) = \int_ {\mathrm{H}} f (x h) \varrho (h) y d \nu (y)
\]

（其中 \(\varphi\in\mathcal{K}(\mathrm{G})\) ，\(y\) 属于 \(\varrho\) 的空间。）

\begin{enumerate}
\def\labelenumi{\arabic{enumi})}
\setcounter{enumi}{11}
\item
  设 G = R，\(\pi\) 为 G 在 \(\mathrm{L}^{2}(\mathbf{R})\) 中的正则表示。证明 \(\chi \leqslant \pi\) 对每个特征标 \(\chi \in \widehat{G}\) 成立。
\item
  设 \(\varrho\) 为 G 在希尔伯特空间 E 中的不可约酉表示，\(\pi\) 为 G 在 F 中的酉表示。设 \(x \in E\) 为范数 1 的向量。假设 \(\varrho \leqslant \pi\) （参见习题~11）。
\end{enumerate}

我们用 \(A \subset \mathcal{C}(G)\) 表示 \(\pi\) 的规范化的对角矩阵系数所成的集合，用 \(\widetilde{A}\) 表示 A 在 \(L^{\infty}(G)\) 中关于弱拓扑的闭凸包。

设 \(x \in E\) 为范数 1 的向量，\(f: g \mapsto \langle x \mid \varrho(g)x \rangle\) 为相伴的对角矩阵系数。

\begin{enumerate}
\def\labelenumi{\alph{enumi})}
\item
  集合 \(\widetilde{\mathrm{A}}\) 是 \(\mathrm{L}^{\infty}(\mathrm{G})\) 中的一个紧凸集，而 \(f\) 是 \(\widetilde{\mathrm{A}}\) 的一个极点。
\item
  函数 f 属于 A 在 \(\widetilde{A}\) 中关于弱拓扑的闭包。
\item
  函数 f 属于 A 在赋予紧收敛拓扑的 \(\mathcal{C}(G)\) 中的闭包。
\item
  G 在 C 中的平凡表示弱包含于 \(\pi\) ，当且仅当对 G 的每个紧子集 C 与每个实数 \(\varepsilon > 0\) ，存在范数为 1 的 \(x \in F\) ，使得
\end{enumerate}

\[
\sup _ {g \in \mathrm{C}} \| \pi (g) x - x \| <   \varepsilon .
\]

\begin{enumerate}
\def\labelenumi{\arabic{enumi})}
\setcounter{enumi}{13}
\tightlist
\item
  设 G 为局部紧群。我们称 G 具有卡日丹性质 ]T[ \(^{(1)}\) ，如果 G 的酉表示 \(\pi\) 具有非零的平凡子表示 \(^{(2)}\) ，当且仅当 G 在 C 上的平凡表示 \(\mathbf{1}_{\text{G}}\) 弱包含于 \(\pi\) 。
\end{enumerate}

\begin{enumerate}
\def\labelenumi{\alph{enumi})}
\item
  设 G 为具有性质 ]T[ 的局部紧群。设 H 为 G 中由 G 的紧子集生成的开子群所成的集合。群 G 是诸子群 H ∈ H 的并。
\item
  设 \(\pi\) 为 G 的这样一个表示，它是 G 在 G/H（\(H \in H\) ）上的拟正则表示的希尔伯特和（参见 V, p. 486 的习题 2）。我们有 \(1_{G} \leqslant \pi\) 。
\item
  群 G 由 G 的一个紧子集生成。特别地，若 G 是离散的，则 G 由 G 的一个有限子集生成。
\end{enumerate}

\begin{enumerate}
\def\labelenumi{\arabic{enumi})}
\setcounter{enumi}{14}
\tightlist
\item
  设 G 为局部紧群，C 为 G 的子集，\(\varepsilon > 0\) 。我们称 \((\mathrm{C}, \varepsilon)\) 是 G 的一个卡日丹对，如果对 G 在希尔伯特空间 E 中的每个酉表示 \(\pi\) ，存在向量 \(x \in E\)
\end{enumerate}

使得

\[
\sup _ {g \in \mathrm{C}} \| \pi (g) x - x \| <   \varepsilon \| x \|,
\]

  便蕴含 \(\pi\) 有一个非零的平凡子表示。我们称 C 为卡日丹集，如果存在 \(\varepsilon > 0\) 使得 \((C, \varepsilon)\) 是一个卡日丹对。

\begin{enumerate}
\def\labelenumi{\alph{enumi})}
\tightlist
\item
  设 \((\mathrm{C}, \varepsilon)\) 为 G 的一个卡日丹对。设 \(\pi\) 为 G 在希尔伯特空间 E 中的酉表示。设 \(x \in E\) 非零且 \(\delta > 0\) ，使得
\end{enumerate}

\[
\sup _ {g \in \mathrm{C}} \| \pi (g) x - x \| <   \delta \varepsilon \| x \|.
\]

则我们有 \(\|p(x)-x\|\leqslant\delta\|x\|\) ，其中 p 是 E 的正交投影算子，其像为 G 不变向量所成的空间。

\begin{enumerate}
\def\labelenumi{\alph{enumi})}
\setcounter{enumi}{1}
\item
  局部紧群 G 具有性质 ]T[，当且仅当存在一个卡日丹对 (C,ε)。（为证性质 ]T[ 蕴含这一断言，用反证法，考虑诸表示 π \(_{C,\varepsilon}\) 的希尔伯特和表示，这些表示由 G 的紧子集 C 与 ε > 0 参数化，使得 π \(_{C,\varepsilon}\) 不含非平凡子表示，但存在范数为 1 的 x \(_{C,\varepsilon}\) 满足 sup \(_{g\in\mathbf{C}}\|\pi(g)x_{\mathrm{C},\varepsilon}-x_{\mathrm{C},\varepsilon}\|<\varepsilon.\)
\item
  设 G 为具有性质 ]T[ 的局部紧群。G 的每个生成 G 的紧子集 C 都是卡日丹集。
\item
  设 G 为具有性质 ]T[ 的局部紧群，C 为 G 的一个紧子集，其内部非空且是卡日丹集。则 C 生成 G。特别地，若 G 是离散的，则卡日丹集恰好是 G 的有限生成子集。（由 C 生成的 G 的子群 H 是开的；考虑 G 在 G/H 上的拟正则表示 \(\pi\) ，并证明 H 的特征函数 \(\varphi\) （H \(\in\) G/H）在 \(\pi\) 的不变向量空间上的投影等于 \(\varphi\) 。）
\end{enumerate}

\begin{enumerate}
\def\labelenumi{\arabic{enumi})}
\setcounter{enumi}{15}
\item
  设 G 为在无穷远处可数的局部紧群，H 为 G 的闭子群。若 G 具有性质 ]T[，且若 G/H 容许一个有界的 G 不变测度，则 H 具有性质 ]T[。（利用 V, p.~492 的习题 20。）
\item
  设 \(n \geqslant 2\) 为整数。我们考虑实李群 \(\mathbf{G} = \mathbf{R}^n \rtimes \mathbf{SL}(n, \mathbf{R})\) ，并用 H 表示 \(\mathbf{R}^n \times \{e\}\) （\(\mathbf{G}\) 的子群），它与 \(\mathbf{R}^n\) 等同。
\end{enumerate}

设 \(\pi\) 为 G 在复希尔伯特空间 E 中的酉表示。假设对 G 的每个紧子集 C 与每个 \(\varepsilon > 0\) ，存在非零的 \(x \in E\) ，使得

\[
\sup _ {g \in \mathrm{C}} \| \pi (g) x - x \| <   \varepsilon \| x \|.
\]

\begin{enumerate}
\def\labelenumi{\alph{enumi})}
\item
  存在 E 中由范数为 1 的向量构成的序列 \((x_{m})\) ，使得函数 \(\varphi_{m}(g)=\langle x_{m}|\pi(g)x_{m}\rangle\) 在赋予紧收敛拓扑的 \(\mathcal{C}(\mathrm{G})\) 中收敛于常值函数 1。（考虑一个其并等于 G 的紧子集序列 \((\mathrm{C}_{m})_{m\in\mathbf{N}}\) ，并取 \(x_{m}\) 使其对 \((\mathrm{C}_{m},1/(m+1))\) 满足上述条件。）
\item
  对每个 \(m \in N\) ，存在总质量为 1 的正测度 \(\mu_{m}\) （在 \(R^{n}\) 上），其傅里叶变换与 \(\varphi_{m}\) 在 H 上的限制等同。
\item
  假设 \(\mu_{m}(\{0\}) = 0\) 对所有 \(m \in N\) 成立。这时把 \(\mu_{m}\) 与 \(R^{n} - \{0\}\) 上总质量为 1 的正测度等同起来，并用 \(\nu_{m}\) 表示 \(\mu_{m}\) 在射影空间 \(\mathbf{P}_{n-1}(\mathbf{R})\) 上的像测度。存在子序列 \((\nu_{n_{k}})_{k \in \mathbf{N}}\) ，它狭收敛于总质量为 1 的正测度 \(\nu\) （在 \(\mathbf{P}_{n-1}(\mathbf{R})\) 上）；测度 \(\nu\) 是 \(\mathbf{SL}(n, \mathbf{R})\) -不变的。由此导出矛盾。
\item
  设 \(m \in N\) 满足 \(\mu_{m}(\{0\}) > 0\) 。设 \(\psi\) 为 \(\varphi_{m}\) 在 H 上的限制。存在 c > 0，使得 \(\psi = c + \widetilde{\psi}\) ，其中 \(\widetilde{\psi}\) 是 \(R^{n}\) 上的正定型函数。
\item
  表示 \(\pi\) 含有 H 不变的向量。（利用习题 2。）我们称 (G,H) 具有相对性质 ]T[。
\end{enumerate}

\begin{enumerate}
\def\labelenumi{\arabic{enumi})}
\setcounter{enumi}{17}
\tightlist
\item
  设 \(\mathrm{G} = \mathbf{SL}(2, \mathbf{R})\) 。用 N 与 P 表示子群
\end{enumerate}

\[
\mathrm{N} = \Big \{\left( \begin{array}{c c} 1 & t \\ 0 & 1 \end{array} \right) | t \in \mathbf {R} \Big \}, \qquad \mathrm{P} = \Big \{\left( \begin{array}{c c} a & t \\ 0 & a ^ {- 1} \end{array} \right) | a \in \mathbf {R} _ {*}, t \in \mathbf {R} \Big \}
\]

（G 的子群）。设 \(\pi\) 为 G 在复希尔伯特空间 E 中的酉表示。假设存在 \(x \in \mathrm{E}\) ，使得 \(\pi(n)x = x\) 对所有 \(n \in \mathrm{N}\) 成立。

\begin{enumerate}
\def\labelenumi{\alph{enumi})}
\tightlist
\item
  齐性空间 \(\mathbf{SL}(2, \mathbf{R}) / \mathrm{N}\) 与 \(\mathbf{R}^2 - \{0\}\) 通过映射 \(g \mapsto ge_1\) 等同，其中 \(e_1\) 是 \(\mathbf{R}^2\) 的典范基的第一个向量；齐性空间 \(\mathbf{SL}(2, \mathbf{R}) / \mathrm{P}\) 与子空间 \((\mathbf{R} - \{0\}) \times \{0\}\) （\(\mathbf{SL}(2, \mathbf{R}) / \mathrm{N}\) 的）等同。
\item
  G 上的每个连续函数 \(f\) ，若满足 \(f(ngn') = f(g)\) （对所有 \((n, n') \in \mathrm{N}^2\) ），则满足 \(f(pgp') = f(g)\) （对所有 \((p, p') \in \mathrm{P}^2\) ）。（把 \(f\) 解释为连续函数 \(\varphi\) （在 \(\mathbf{R}^2 - \{0\}\) 上、在 N 的自然作用下不变），并验证它在 \(\mathbf{SL}(2, \mathbf{R}) / \mathrm{P}\) 上是常数。）
\end{enumerate}

\begin{enumerate}
\def\labelenumi{\alph{enumi})}
\setcounter{enumi}{2}
\tightlist
\item
  有 \(\pi(g)x = x\) 对所有 \(g \in G\) 成立。
\end{enumerate}

\begin{enumerate}
\def\labelenumi{\arabic{enumi})}
\setcounter{enumi}{18}
\tightlist
\item
  设 \(n \geqslant 2\) 为整数。考虑实李群 \(\mathrm{G} = \mathbf{SL}(n, \mathbf{R})\) 。用 N 与 H 表示 G 中元素分别为分块形式
\end{enumerate}

\[
\left( \begin{array}{c c} \operatorname{Id} _ {n - 1} & x \\ 0 & 1 \end{array} \right), \quad (x \in \mathbf {R} ^ {n - 1}), \qquad \left( \begin{array}{c c} g & 0 \\ 0 & 1 \end{array} \right), \quad (g \in \mathbf {S L} (n - 1, \mathbf {R}))
\]

的闭子群。群 N 与 \(R^{n-1}\) 等同，群 H 与 \(\mathbf{SL}(n-1,\mathbf{R})\) 等同。

\begin{enumerate}
\def\labelenumi{\alph{enumi})}
\item
  设 \(\pi\) 为 G 在复希尔伯特空间 E 中的酉表示。E 的元素 \(x\) 在 G 下不变，当且仅当它在 N 下不变。
\item
  若 \(n \geqslant 3\) ，则群 G 具有性质 ]T[。（把上一问题的结果与上一习题结合起来。）
\item
  若 \(n \geqslant 3\) ，则群 \(\widetilde{G} = R^{n} \rtimes G\) 具有性质 ]T[。（设 \(\pi\) 为满足 \(1_{\widetilde{G}} \leqslant \pi\) 的酉表示；由上一问题，存在非零向量 \(x \neq 0\) ，它被 \(\widetilde{G}\) 的子群 G 保持不变；证明相应的矩阵系数是常数。）
\item
  若 \(n \geqslant 3\) ，则群 \(\mathbf{SL}(n, \mathbf{Z})\) 具有性质 ]T[。（利用 INT, VII, p.~116 的习题 7。）
\item
  设 \(n \geqslant 3\) 。群 \(\mathbf{SL}(n, \mathbf{Z})\) 是有限生成的。设 H 为 \(\mathbf{SL}(n, \mathbf{Z})\) 的有限指标正规子群所成的族；它是无限的。对 \(\mathbf{SL}(n, \mathbf{Z})\) 的每个有限生成子集 S，凯莱图族 \((\mathcal{G}(\mathbf{SL}(n, \mathbf{Z})/\mathrm{H}, \mathrm{SH}/\mathrm{H}))_{\mathrm{H} \in \mathcal{H}}\) 是一个扩展图族（参见 TA, 2, p.~221, 习题 7 及 IV, p.~323 的习题 16）。
\end{enumerate}

\begin{enumerate}
\def\labelenumi{\arabic{enumi})}
\setcounter{enumi}{19}
\tightlist
\item
  设 G 为局部紧群，\(\mu\) 为 G 上的左哈尔测度。我们用 \(\Delta\) 表示 G 的模。对 G 上的任意函数 f，我们将记 \(\widetilde{f} = \Delta f^{*}\) 。
\end{enumerate}

G 上的测度 \(\nu\) 称为正定型的，如果 \(\langle \nu, \tilde{f} * f \rangle \geqslant 0\) 对每个函数 \(f \in \mathcal{K}(\mathrm{G})\) 成立。G 上的局部可积函数 \(\varphi\) 称为正定型的，如果 \(\langle \varphi, f^* * f \rangle \geqslant 0\) 对每个函数 \(f \in \mathcal{K}(\mathrm{G})\) 成立。

\begin{enumerate}
\def\labelenumi{\alph{enumi})}
\item
  假设 \(\nu\) 有界。测度 \(\nu\) 是正定型的，当且仅当 \(\gamma(\nu) \geqslant 0\) 。
\item
  设 \(\varphi\in\mathcal{K}(\mathrm{G})\) 。函数 \(\varphi\) 作为局部可积函数是正定型的，当且仅当它按照 V, p.~442 的定义 4 是正定型的。
\item
  设 \(\nu\) 为 G 上的一个正定型测度。公式
\end{enumerate}

\[
(f \mid g) _ {\nu} = \langle \nu , \widetilde {f} * g \rangle
\]

在 \(\mathcal{K}(\mathrm{G})\) 上定义了一个正埃尔米特型。我们用 \(\mathrm{E}_{\nu}\) 表示赋予了该埃尔米特型的 \(\mathcal{K}(\mathrm{G})\) 的分离完备希尔伯特空间。

\begin{enumerate}
\def\labelenumi{\alph{enumi})}
\setcounter{enumi}{3}
\tightlist
\item
  对 \(f \in \mathcal{K}(\mathrm{G})\) ，函数 \(\sigma(g)f \colon g \mapsto (\varepsilon_g * f)\Delta^{1/2}\) 属于 \(\mathcal{K}(\mathrm{G})\) ，且 \(\| \sigma(g)f \|_{\nu} = \|f\|_{\nu}\) 。映射 \(g \mapsto \sigma(g)\) 延拓为 G 在 \(\mathrm{E}_{\nu}\) 中的一个酉表示。
\end{enumerate}

\begin{enumerate}
\def\labelenumi{\alph{enumi})}
\setcounter{enumi}{4}
\tightlist
\item
  设 U 为 G 中 e 的开邻域，\(\nu\) 为 G 上的正定型测度，使得存在实数 \(c \geqslant 0\) 满足
\end{enumerate}

\[
\langle \nu , f * \widetilde {f} \rangle \leqslant c \left(\int_ {\mathrm{G}} f d \mu\right) ^ {2}
\]

  其中 \(f \in \mathcal{K}(\mathrm{G})\) 为正的且在 U 外为零。存在 G 上的连续正定型函数 \(\varphi\) ，使得 \(\nu = \varphi \Delta^{-1/2} \cdot \mu\) 。（设 \(j\) 为从 \(\mathcal{K}(\mathrm{G})\) 到 \(\mathrm{E}_{\nu}\) 的典范态射；若 \(\mathfrak{B}\) 是 \(e\) 的邻域滤子的一组基（其元素都含于 U），而 \(f_{\mathrm{V}} \in \mathcal{K}_{+}(\mathrm{G})\) 是支集含于 \(V \in B\) 、积分为 1 的函数，证明 \(j(f_{\mathrm{V}})\) 沿 B 狭收敛于元素 \(y_{0}\) （属于 \(E_{\nu}\) ），使得

\[
(j (f) \mid y _ {0}) = \int_ {\mathrm{G}} f ^ {*} d \mu
\]

对 \(f \in \mathcal{K}(\mathrm{G})\) 成立；令 \(\varphi(g) = (y_0 \mid \sigma(g)y_0)\) ，验证 \(\varphi\) 具有所需的性质。

\begin{enumerate}
\def\labelenumi{\alph{enumi})}
\setcounter{enumi}{5}
\tightlist
\item
  设 \(\varphi_{1}\) 为 G 上的连续正定函数，\(\varphi_{2}\) 为局部可积的正定型函数。若 \(\varphi_{1}-\varphi_{2}\) 是正定型的，则 \(\varphi_{2}\) 局部几乎处处等于一个连续的正定型函数。（验证可以应用上一问题。）
\end{enumerate}

\begin{enumerate}
\def\labelenumi{\arabic{enumi})}
\setcounter{enumi}{20}
\tightlist
\item
  我们沿用上一习题的记号。
\end{enumerate}

\begin{enumerate}
\def\labelenumi{\alph{enumi})}
\item
  设 \(f \in \mathcal{L}^{2}(G)\) 作为局部可积函数是正定型的。\(L^{2}(G)\) 上以 \(\mathcal{K}(G)\) 为定义域、由 \(\varphi \mapsto \gamma(\varphi)f\) 定义的部分算子是正的。我们用 \(\varrho(f)\) 表示 f 的弗里德里希斯延拓（IV, p.~295 的例 8）；它是 \(L^{2}(G)\) 上的一个正自伴部分算子。
\item
  设 \(f \in \mathcal{L}^2(\mathrm{G})\) 。我们称 \(f\) 是清醒的，如果存在实数 \(c \geqslant 0\) ，使得 \(\| \boldsymbol{\gamma}(\varphi)f \| \leqslant c \| \varphi \|\) 对所有 \(\varphi \in \mathcal{K}(\mathrm{G}) \subset \mathcal{L}^2(\mathrm{G})\) 成立。这时我们用 \(\varrho(f)\) 表示 \(\mathrm{L}^2(\mathrm{G})\) 中延拓 \(\varphi \mapsto \boldsymbol{\gamma}(\varphi)f\) 的自同态。若 \(f\) 是正定型的，这一定义与上一个定义一致，并且这一条件此时等价于说 \(\varrho(f)\) 是正的。
\item
  设 \(f_{1}\) 与 \(f_{2}\) 为 \(\mathcal{L}^{2}(G)\) 中的清醒函数。若 \(\varrho(f_{1}) = \varrho(f_{2})\) ，则 \(f_{1}\) 与 \(f_{2}\) 在 \(L^{2}(G)\) 中相等。
\end{enumerate}

\begin{enumerate}
\def\labelenumi{\alph{enumi})}
\setcounter{enumi}{3}
\tightlist
\item
  设 \(f \in \mathcal{L}^2(\mathrm{G})\) 为正定型的。对所有 \(g \in \mathrm{G}\) ，我们有 \(\varrho(f) \circ \boldsymbol{\gamma}(g) = \boldsymbol{\gamma}(g) \circ \varrho(f)\) 。此外，对每个 \(h \in \operatorname{dom}(\varrho(f))\) ，我们有 \(\varrho(f)h = h * f\) 。
\end{enumerate}

\begin{enumerate}
\def\labelenumi{\alph{enumi})}
\setcounter{enumi}{4}
\tightlist
\item
  设 \(f_{1}\) 与 \(f_{2}\) 为 \(\mathcal{L}^{2}(\mathrm{G})\) 中的清醒元素，使得 \(\varrho(f_{1})\) 与 \(\varrho(f_{2})\) 交换。函数 \(f_{1}*f_{2}\) 是 G 上的连续正定函数；它是清醒的，并且我们有
\end{enumerate}

\[
\varrho (f _ {1} * f _ {2}) = \varrho (f _ {1}) \varrho (f _ {2}), \quad \langle f _ {1} | f _ {2} \rangle \geqslant 0.
\]

\begin{enumerate}
\def\labelenumi{\alph{enumi})}
\setcounter{enumi}{5}
\tightlist
\item
  在上一问题的记号与假设下，若 \(f_{2} - f_{1}\) 是正定型的，则
\end{enumerate}

\[
\left\| f _ {2} - f _ {1} \right\| ^ {2} \leqslant \left\| f _ {2} \right\| ^ {2} - \left\| f _ {1} \right\| ^ {2}.
\]

\begin{enumerate}
\def\labelenumi{\alph{enumi})}
\setcounter{enumi}{6}
\tightlist
\item
  设 \((f_n)_{n\in \mathbf{N}}\) 为 \(\mathcal{L}^2 (\mathrm{G})\) 中两两交换的清醒正定型元素构成的序列，且 \(f_{n + 1} - f_n\) 对所有 \(n\in \mathbf{N}\) 是正定型的。若序列 \((\| f\| _n)\) 有界，则序列 \((f_n)\) 在 \(\mathcal{L}^2 (\mathrm{G})\) 中收敛。
\end{enumerate}

\begin{enumerate}
\def\labelenumi{\arabic{enumi})}
\setcounter{enumi}{21}
\tightlist
\item
  我们沿用上一习题的记号。
\end{enumerate}

设 \(f\) 为属于 \(\mathcal{L}^2 (\mathrm{G})\) 的连续正定型函数。

\begin{enumerate}
\def\labelenumi{\alph{enumi})}
\item
  存在 [0,1] 上多项式函数的递增序列 \((p_{n})_{n\in\mathbf{N}}\) ，它在 [0,1] 上一致收敛于函数 \(t\mapsto\sqrt{t}\) 。
\item
  假设 \(f\) 是清醒的，且 \(0 \leqslant f \leqslant 1\) （在 \(\mathcal{L}(\mathrm{L}^2(\mathrm{G}))\) 中）。存在序列 \((h_n)\) （在 \(\mathcal{L}^2(\mathrm{G})\) 中），满足下列条件：
\end{enumerate}

\begin{enumerate}
\def\labelenumi{(\roman{enumi})}
\item
  每个 \(h_n\) 都是清醒的且是正定型的；
\item
  函数 \(h_{n+1} - h_n\) 对所有 \(n \in \mathbf{N}\) 是正定型的；
\item
  自同态 \(\varrho(h_n)\) 两两交换；
\item
  有 \(\varrho(h_n)^2 \leqslant \varrho(f)\) （对所有 \(n\) ）。
\end{enumerate}

\begin{enumerate}
\def\labelenumi{\alph{enumi})}
\setcounter{enumi}{2}
\tightlist
\item
  序列 \((h_n)_{n\in \mathbf{N}}\) 在 \(\mathcal{L}^2 (\mathrm{G})\) 中收敛；若 \(h\) 是 \((h_n)\) 的一个极限，则 \(h\) 是正定型且清醒的，并满足 \(\varrho (h)^2 = \varrho (f)\) ；我们有 \(f = h*h\) 。
\end{enumerate}

\begin{enumerate}
\def\labelenumi{\alph{enumi})}
\setcounter{enumi}{3}
\tightlist
\item
  我们不再假设 \(f\) 是清醒的。对 \(t \in \mathbf{R}\) ，用 \(p_t\) 表示部分自伴算子 \(\varrho(f)\) 的由 \([- \infty, t]\) 定义的谱投影。我们有 \(p_t \circ \boldsymbol{\gamma}(g) = \boldsymbol{\gamma}(g) \circ p_t\) 对所有 \(t \in \mathbf{R}\) 与 \(g \in G\) 成立。
\end{enumerate}

\begin{enumerate}
\def\labelenumi{\alph{enumi})}
\setcounter{enumi}{4}
\tightlist
\item
  对 \(t \in \mathbf{R}\) ，令 \(f_t = p_t(f)\) ；它是 \(\mathrm{L}^2(\mathrm{G})\) 中某个连续正定的清醒函数（属于 \(\mathcal{L}^2(\mathrm{G})\) ）的类。（验证 \(\gamma(\varphi)f_t = \varrho(f)p_t(g)\) 对 \(g \in \mathcal{H}(\mathrm{G})\) 成立，由此推出 \(\varrho(f_t) = \varrho(f) \circ p_t\) ，再应用习题 20 的结果。）
\end{enumerate}

\begin{enumerate}
\def\labelenumi{\alph{enumi})}
\setcounter{enumi}{5}
\tightlist
\item
  存在 \(h \in \mathcal{L}^2(\mathrm{G})\) ，使得 \(f = h * h\) 。（考虑 \(f_n = p_n(f)\) （\(n \in \mathbf{N}\) ），并应用部分 \(c\) 。）
\end{enumerate}

\begin{enumerate}
\def\labelenumi{\arabic{enumi})}
\setcounter{enumi}{22}
\tightlist
\item
  我们记 \(G = R_{+}^{*}\) ；它是一个局部紧交换群，且测度 \(\mu = dx/x\) 是 G 上的哈尔测度。群 G 与群 iR 通过映射 \((x, it) \mapsto x^{it}\) 处于对偶；\(\mu\) 的对偶测度是 iR 上的勒贝格测度（参见 II, p. 236 的推论 4）。我们用 F 表示 G 的傅里叶变换，用 \(\overline{F}\) 表示 \(\widehat{G}\) 的傅里叶余变换。
\end{enumerate}

我们用 X 表示 G 上的恒等函数。

\begin{enumerate}
\def\labelenumi{\alph{enumi})}
\item
  设 \(\nu\) 为 G 上的正测度。我们用 \(\mathrm{I}_{\nu}\) 表示这样的 \(\sigma \in \mathbf{R}\) 所成的集合：函数 \(\mathrm{X}^{\sigma}\) 是 \(\nu\) -可积的；用 \(\mathrm{B}_{\nu}\) 表示这样的 \(s \in \mathbf{R}\) 所成的集合：\(\mathcal{R}(s) \in \mathrm{I}_{\nu}\) 。集合 \(\mathrm{I}_{\nu}\) 是 \(\mathbf{R}\) 的一个区间，而函数 \(s \mapsto \nu(\mathrm{X}^s)\) 在 \(\mathrm{B}_{\nu}\) 的内部有定义且全纯。
\item
  设 \(f\) 为 G 上的局部可积函数。我们记 \(\mathrm{I}_{f} = \mathrm{I}_{f\cdot \mu}\) 与 \(\mathrm{B}_f = \mathrm{B}_{f\cdot \mu}\) 。对每个 \(s \in \mathrm{B}_f\) ，我们令
\end{enumerate}

\[
\widehat {f} (s) = \int_ {\mathrm{G}} f (x) x ^ {s} d \mu (x).
\]

若存在 \(\delta > 0\) 与 \(c \geqslant 0\) ，使得 \(|\widehat{f}(s)| \leqslant c(1 + |s|)^{-1 - \delta}\) 对 \(s \in \mathrm{B}_f\) 成立，则对所有 \(\sigma\) （属于 \(\mathrm{B}_f\) 的内部），我们有

\[
f (x) = \frac {1}{2 \pi} \int_ {\mathbf {R}} \widehat {f} (\sigma + i t) x ^ {- \sigma - i t} d t.
\]

\begin{enumerate}
\def\labelenumi{\alph{enumi})}
\setcounter{enumi}{2}
\tightlist
\item
  设 \(f \in \mathcal{L}^{2}(\mathrm{G})\) 。假设 0 属于 \(I_{f}\) 的内部。则我们有 \(\overline{\mathcal{F}} f(it) = \widehat{f}(it)\) 对几乎所有 \(t \in \mathbf{R}\) 成立。
\end{enumerate}

\begin{enumerate}
\def\labelenumi{\alph{enumi})}
\setcounter{enumi}{3}
\item
  设 \(f\) 与 \(g\) 为 \(\mathbf{G}\) 上的可测函数。设 \(\sigma\) 为实数，使得 \(\mathrm{X}^{-2\sigma}f\in \mathrm{L}^2 (\mathrm{G})\) 且 \(\mathrm{X}^{2\sigma}g\in \mathrm{L}^2 (\mathrm{G})\) 。则我们有

\[
\int_ {\mathrm{G}} f g d \mu = \frac {1}{2 \pi} \int_ {\mathbf {R}} \widehat {f} (- \sigma - i t) \widehat {g} (\sigma + i t) d t.
\]

\end{enumerate}

\begin{enumerate}
\def\labelenumi{\alph{enumi})}
\setcounter{enumi}{4}
\tightlist
\item
  设 \(f\) 与 \(g\) 为 \(G\) 上的可测函数。假设
\end{enumerate}

\begin{enumerate}
\def\labelenumi{(\roman{enumi})}
\item
  存在 \(c \in \mathbf{R}\) ，使得 \(|f(x)| \leqslant cx / (1 + x)\) 对所有 \(x \in \mathbf{G}\) 成立；
\item
  有 ]0,1[ ⊂ I \(_{g}\) ，且对每个满足 0 < σ < 1 的 σ，函数 X \(^{\sigma}\) g 有界；
\item
  存在 \(c' \in \mathbf{R}\) ，使得 \(|\widehat{g}(s)| \leqslant c'(1 + |s|)^{-2}\) 对所有 \(s \in \mathbf{C}\) （满足 \(\mathcal{R}(s) \in ]0, 1[\) ）成立。
\end{enumerate}

则我们有

\[
\int_ {\mathrm{G}} f g d \mu = \frac {1}{2 \pi} \int_ {\mathbf {R}} \widehat {f} (- \sigma - i t) \widehat {g} (\sigma + i t) d t
\]

对所有 \(\sigma\in\mathbf{R}\) （满足 \(0<\sigma<1\) ）成立。

\begin{enumerate}
\def\labelenumi{\alph{enumi})}
\setcounter{enumi}{5}
\tightlist
\item
  设 \(\sigma \in \mathbf{R}\) ，\(f \in \mathcal{C}(\sigma + i\mathbf{R})\) 。G 上存在正测度 \(\nu\) ，使得 \(\sigma \in \mathrm{I}_{\nu}\) 且 \(\widehat{\nu}(\sigma + it) = f(\sigma + it)\) 对所有 \(t \in \mathbf{R}\) 成立，当且仅当对每个整数 \(n \geqslant 1\) 、复数的每个族 \((t_i)_{0 \leqslant i \leqslant n}\) 与任意族 \((z_i)_{0 \leqslant i \leqslant n}\) （\(\mathcal{R}(z_i) = \sigma\) 对所有 \(i\) ），我们有
\end{enumerate}

\[
\sum_ {i = 0} ^ {n} \sum_ {j = 0} ^ {n} \bar {t} _ {i} t _ {j} f \left(\frac {1}{2} (\overline {{z}} _ {i} + z _ {j})\right) \geqslant 0.
\]

\exercisesection{}

\begin{enumerate}
\def\labelenumi{\arabic{enumi})}
\tightlist
\item
  设 G 为紧拓扑群，B 为 G 的不可约酉表示的特征标所成的族。
\end{enumerate}

\begin{enumerate}
\def\labelenumi{\alph{enumi})}
\item
  偶 \((\mathrm{R}(\mathrm{G}),\mathcal{B})\) 是一个自然的含幺基代数（参见 III, p.~130 的习题 6）。
\item
  假设 G 有限。代数 R(G) 这时是有限秩的；计算 \(p_f(\pi)\) （对每个 \(\pi \in \widehat{G}\) ），并辨识与上述引文问题 j) 中定义的元素 r₀ 相对应的 G 的表示。
\end{enumerate}

\begin{enumerate}
\def\labelenumi{\arabic{enumi})}
\setcounter{enumi}{1}
\tightlist
\item
  设 \(n \geqslant 1\) ，\(p \in [1, +\infty]\) 。我们给 \(\mathbf{R}^n\) 赋予范数
\end{enumerate}

\[
\| x \| = \left(\sum_ {i = 1} ^ {n} | x _ {i} | ^ {p}\right) ^ {1 / p} \text {   if   } p \neq + \infty ,
\]

而 \(\|x\| = \sup_{i}|x_{i}|\) （若 \(p = +\infty\) ）。设 G 为赋有该范数的度量空间 \(R^{n}\) 的等距群。

\begin{enumerate}
\def\labelenumi{\alph{enumi})}
\item
  设 \(\sigma\) 为 \(\{1,\ldots,n\}\) 的一个置换，\(\varepsilon\in\{-1,1\}^{n}\) 。设 \((e_{i})\) 为 \(\mathbf{R}^{n}\) 的典范基。\(\mathbf{R}^{n}\) 中使 \(e_{i}\mapsto\varepsilon_{i}e_{\sigma(i)}\) 的唯一自同态属于 G。这种形式的等距所成的集合是 G 的一个子群 W；它是 \(\mathfrak{S}_{n}\) 被 \((\mathbf{Z}/2\mathbf{Z})^{n}\) 的一个扩张。
\item
  群 G 是紧的。
\item
  群 G 在 \(R^{n}\) 上的作用是不可约的；群 G 含于二次型 \(\sum x_{i}^{2}\) 的正交群。
\item
  若 \(p \neq 2\) ，则 G = W。
\end{enumerate}

\begin{enumerate}
\def\labelenumi{\arabic{enumi})}
\setcounter{enumi}{2}
\tightlist
\item
  我们回顾，称群 \(\Gamma\) 是剩余有限的，如果 \(\Gamma\) 的有限指标子群的交约化为 \(\{e\}\) （A, I, p.~129, 习题 5, b)）。
\end{enumerate}

对群 \(\Gamma\) ，我们用 \(\widehat{\Gamma}\) 表示射影极限 \(\widehat{\Gamma} = \varprojlim \Gamma / H\) ，其中 \(H\) 取遍 \(\Gamma\) 的按包含排序的有限指标正规子群所成的集合。我们称 \(\widehat{\Gamma}\) 是 \(\Gamma\) 的射有限完备化。

\begin{enumerate}
\def\labelenumi{\alph{enumi})}
\item
  群 \(\Gamma\) 是剩余有限的，当且仅当对每个异于单位元的 \(g \in \Gamma\) ，存在有限群 F 与同态 \(\varphi: \Gamma \to F\) ，使得 \(\varphi(g) \neq e\) 。
\item
  证明 \(\Gamma\) 是剩余有限的，当且仅当从 \(\Gamma\) 到 \(\widehat{\Gamma}\) 的典范同态是单射。
\item
  设 K 为域。\(\mathbf{GL}(n,\mathrm{K})\) 的每个有限生成子群（参见 A, I, p.~145, 习题 16）都是剩余有限的（“马尔采夫定理”）。（证明存在整数 \(m\geqslant 0\) 、一个等于 \(\mathbf{Z}\) 或某个有限域的环 A，以及该子群到 \(\mathbf{GL}(m,\mathrm{R})\) 的一个单射同态，其中 R 形如 \(\mathrm{A}[\mathrm{X}_1,\dots ,\mathrm{X}_r,1 / f]\) 。）
\item
  设 \(\Gamma\) 为有限生成的离散群。假设存在从 \(\Gamma\) 到某个紧拓扑群的单射同态。对每个异于单位元的 \(\gamma \in \Gamma\) ，存在 \(\Gamma\) 的一个有限维不可约酉表示，使得 \(\gamma\) 在 \(\pi\) 下的像不是恒等映射。
\item
  有限生成群是剩余有限的，当且仅当它同构于某个紧群的一个子群。
\end{enumerate}

¶ 4) 设 G 为群，A 与 B 为 G 的子群。设 \(\varphi\colon A\to B\) 是一个同构。

设 \(r'\) 为自由群 F(G) 中的一个族，使得 G 同构于群 F(G, \(r'\) ) （A, I, p.~87, n°6）。我们称 G 依 \(\varphi\) 而得的 HNN 扩张为群 \(G^{\varphi} = F(X, r)\) ，其中 X 是集合 G、A、B 与一个元素 t 之和，而 r 是 \(r'\) 与关系子 \(t^{-1}at = \varphi(a)\) （对所有 \(a \in A\) ）的和族。

我们用 \(j\colon G\to G^{\varphi}\) 表示唯一的同态，使得 \(j(g)\) 是 \(g\) 在 \(G^{\varphi}\) 中所代表的元素。

\begin{enumerate}
\def\labelenumi{\alph{enumi})}
\tightlist
\item
  设 \(k \geqslant 1\) 为整数，\((g_{0}, \varepsilon_{1}, g_{1}, \ldots, \varepsilon_{k}, g_{k})\) 为一个有限族，使得对每个 i 有 \(g_{i} \in G\) 且 \(\varepsilon_{i} \in \{-1, 1\}\) 。假设不存在满足 \(1 \leqslant j \leqslant k - 1\) 的整数 j 使得
\end{enumerate}

\[
(\varepsilon_ {j}, \varepsilon_ {j + 1}) = (- 1, 1), \quad g _ {i} \in A
\]

或

\[
(\varepsilon_ {j}, \varepsilon_ {j + 1}) = (1, - 1), \quad g _ {i} \in \mathrm{B}.
\]

则元素

\[
j (g _ {0}) t ^ {\varepsilon_ {1}} j (g _ {1}) \dots t ^ {\varepsilon_ {k}} j (g _ {k})
\]

（G \(^{φ}\) 的这个元素）不是单位元（“布里顿引理”。）

\begin{enumerate}
\def\labelenumi{\alph{enumi})}
\setcounter{enumi}{1}
\item
  同态 j 是单射。
\item
  设 \(a\) 与 \(b\) 为不同的符号。设 G 为群 \(\mathrm{F}(\{a,b\}, a^{-1}b^{2}ab^{-3})\) 。记 \(b_{1} = a^{-1}ba\) 。设 H 为 G 的一个有限指标正规子群，\(\pi: \mathrm{G} \to \mathrm{G}/\mathrm{H}\) 为典范投影。有 \([b_{1}, b] \in \mathrm{Ker}(\pi)\) 。
\item
  群 G 不容许忠实的有限维表示。（证明 G 不是剩余有限的。）
\end{enumerate}

\begin{enumerate}
\def\labelenumi{\arabic{enumi})}
\setcounter{enumi}{4}
\tightlist
\item
  设 G 为紧群，\(\varrho\) 为 G 的一个忠实的有限维酉表示。设 \(\pi \in \widehat{G}\) 。
\end{enumerate}

\begin{enumerate}
\def\labelenumi{\alph{enumi})}
\tightlist
\item
  存在整数 a 与 b，使得 \(\pi\) 同构于 \(\varrho^{\otimes a} \otimes \overline{\varrho}^{\otimes b}\) 的一个子表示。（例如可以利用魏尔斯特拉斯–斯通定理。）
\end{enumerate}

我们用 \(w(\pi)\) 表示 \(a + b\) 当 \((a, b)\) 取遍满足这一性质的整数组所成的集合时的最小值。

\begin{enumerate}
\def\labelenumi{\alph{enumi})}
\setcounter{enumi}{1}
\tightlist
\item
  设 \(\Lambda\) 为 \(G^{0}\) 的权格，q 为 \(\Lambda\otimes \mathbf{R}\) 上的一个范数。存在实数 C 与 D，使得对每个表示 \(\pi\in\widehat{G}\) （其在 \(G^{0}\) 上的限制具有最高权 \((\lambda_{1},\ldots,\lambda_{k})\) ），我们有
\end{enumerate}

\[
w (\pi) \leqslant \mathrm{C} \sup (q (\lambda_ {1}), \dots , q (\lambda_ {k})) + \mathrm{D}.
\]

（先考虑连通半单群的情形，并对 G 的适当选取的有限表示族应用 a)。）

\begin{enumerate}
\def\labelenumi{\arabic{enumi})}
\setcounter{enumi}{5}
\tightlist
\item
  设 G 为紧群，\(\varrho\) 为 G 的一个不可约表示。存在中心函数 \(f \in \Theta(G^{\sharp})\) ，使得
\end{enumerate}

\[
f = \sum_ {\pi \in \widehat {\mathrm{G}}} c (\pi) \chi_ {\pi}
\]

其中整数 \(c(\pi)\) 满足

\begin{enumerate}
\def\labelenumi{(\roman{enumi})}
\item
  有 \(c(\pi) \geqslant 0\) （对所有 \(\pi \in \widehat{G}\) ）；
\item
  \(c(1) \leqslant c(\varrho)\) ，且除 \(\varrho\) 是一维的并满足 \(\varrho^{2} = 1\) 的可能情形外，该不等式是严格的；
\item
  \(\mathcal{R}(f) \geqslant 0.\)
\end{enumerate}

¶ 7) 对每个整数 \(n \geqslant 1\) ，我们用 \(\mu_n\) 表示 \(\mathbf{C}\) 上群 \(\mathbf{U}(\mathbf{C}^n) \subset \mathrm{GL}(\mathbf{C}^n)\) 的规范化哈尔测度在迹映射 \(g \mapsto \mathrm{Tr}(g)\) 下的像；它是总质量为 1 的紧支集测度。

\begin{enumerate}
\def\labelenumi{\alph{enumi})}
\tightlist
\item
  设 \(a\) 与 \(b\) 为自然数，满足 \(a + b \leqslant n\) 。我们有
\end{enumerate}

\[
\int_ {\mathbf {C}} z ^ {a} \overline {{z}} ^ {b} d \mu_ {n} (z) = \left\{ \begin{array}{l l} 0 & \text{若 } a \neq b \\ a! & \text{若 } a = b. \end{array} \right.
\]

\begin{enumerate}
\def\labelenumi{\alph{enumi})}
\setcounter{enumi}{1}
\tightlist
\item
  测度序列 \((\mu_{n})_{n\geqslant1}\) 狭收敛于实向量空间 C 上协方差为 \(\mathrm{Q}(z)=|z|^{2}\) 的高斯测度。
\end{enumerate}

\begin{enumerate}
\def\labelenumi{\arabic{enumi})}
\setcounter{enumi}{7}
\item
  设 G 为紧群，\(\varrho_{1}\) 与 \(\varrho_{2}\) 为 G 的不可约酉表示。设 \(n \geqslant 2\) 为整数。若 \(S^{n}\varrho_{1}\) 同构于 \(S^{n}\varrho_{2}\) ，则存在 G 的一个有限指标子群 H，使得 \(\varrho_{1}\) 在 H 上的限制同构于 \(\varrho_{2}\) 在 H 上的限制。
\item
  设 \(k\) 为含 5 个元素的有限域，\(\mathrm{G} = \mathbf{S}\mathbf{L}(2,k)\) 。
\end{enumerate}

\begin{enumerate}
\def\labelenumi{\alph{enumi})}
\tightlist
\item
  群 G 的共轭类由 e、-e 以及七个元素
\end{enumerate}

\[
g _ {3} = \left( \begin{array}{c c} 1 & 1 \\ 0 & 1 \end{array} \right), \qquad g _ {4} = \left( \begin{array}{c c} 1 & 2 \\ 0 & 1 \end{array} \right), \qquad g _ {5} = \left( \begin{array}{c c} - 1 & 1 \\ 0 & - 1 \end{array} \right), \qquad g _ {6} = \left( \begin{array}{c c} - 1 & 2 \\ 0 & - 1 \end{array} \right)
\]

\[
g _ {7} = \left( \begin{array}{c c} 2 & 0 \\ 0 & - 2 \end{array} \right), \qquad g _ {8} = \left( \begin{array}{c c} 0 & 1 \\ - 1 & 1 \end{array} \right), \qquad g _ {9} = \left( \begin{array}{c c} 0 & 1 \\ - 1 & - 1 \end{array} \right).
\]

代表。确定每个共轭类的基数。

\begin{enumerate}
\def\labelenumi{\alph{enumi})}
\setcounter{enumi}{1}
\tightlist
\item
  设 \(\alpha = \frac{1}{2}(1 + \sqrt{5})\) 。存在 G 的一个二维忠实不可约酉表示 \(\pi\) ，其特征标 \(\chi\) 满足 \(\chi(-e) = -2\) 以及
\end{enumerate}

\[
\begin{array}{c} \chi (g _ {4}) = \alpha , \quad \chi (g _ {6}) = - \alpha , \quad \chi (g _ {3}) = - \alpha - 1, \quad \chi (g _ {5}) = \alpha + 1 \\ \chi (g _ {7}) = 0, \quad \chi (g _ {8}) = 1, \quad \chi (g _ {9}) = - 1. \end{array}
\]

（可以考虑 \(\mathbf{GL}(2,k)\) 的四维不可约表示在 G 上的限制，参见 A, VIII, p.~422, 习题 27。）

\begin{enumerate}
\def\labelenumi{\alph{enumi})}
\setcounter{enumi}{2}
\tightlist
\item
  有 \(\mathrm{M}_4(\pi) = 2\) 。
\end{enumerate}

\begin{enumerate}
\def\labelenumi{\arabic{enumi})}
\setcounter{enumi}{9}
\tightlist
\item
  \begin{enumerate}
  \def\labelenumii{\alph{enumii})}
  \tightlist
  \item
    存在一个正交型的不可约表示 \(\pi\) （交错群 \(\mathfrak{A}_5\) 的），使得 \(\mathrm{M}_4(\pi) = 3\) 。
  \end{enumerate}
\end{enumerate}

\begin{enumerate}
\def\labelenumi{\alph{enumi})}
\setcounter{enumi}{1}
\tightlist
\item
  设 R 为 \(E_{6}\) 型（相应地 \(E_{7}\) 型、\(E_{8}\) 型）的根系，W 为 R 的外尔群。设 \(\pi\) 为 W 的几何表示（LIE, V, § 4, 第 3 目）。它是不可约的、正交型的，且满足 \(M_{4}(\pi)=3\) 。
\end{enumerate}

（可以利用诸如 GAP 或 MAGMA 的软件来确定 W 的共轭类及其在 \(\pi\) 下的像的列表。）

\begin{enumerate}
\def\labelenumi{\arabic{enumi})}
\setcounter{enumi}{10}
\item
  设 G 为紧群，\(\pi\) 为 G 在希尔伯特空间 E 中的有限维酉表示。若 \(\mathrm{M}_{4}(\pi) \leqslant 5\) ，则表示 \(\pi\) 是不可约的。
\item
  设 G 为连通紧半单李群，其复化李代数同构于单李代数 \(e_{6}\) 。
\end{enumerate}

\begin{enumerate}
\def\labelenumi{\alph{enumi})}
\item
  存在 G 的一个 27 维忠实不可约表示 \(\varrho\) ；这样的表示是复型的。
\item
  表示 \(\varrho\) 的四阶矩等于 3。
\end{enumerate}

\begin{enumerate}
\def\labelenumi{\arabic{enumi})}
\setcounter{enumi}{12}
\tightlist
\item
  设 G 为连通紧半单李群，g 为其李代数。假设存在 G 的一个 27 维忠实不可约表示 \(\varrho\) 。我们假设 \(\varrho\) 的四阶矩等于 3，且表示 \(\varrho\) 是复型的。
\end{enumerate}

\begin{enumerate}
\def\labelenumi{\alph{enumi})}
\item
  李代数 g 是单的。
\item
  李代数 g 同构于 \(\mathfrak{e}_{6}\) 。
\end{enumerate}

\begin{enumerate}
\def\labelenumi{\arabic{enumi})}
\setcounter{enumi}{13}
\item
  设 G 为连通紧半单李群。假设存在 G 的一个 27 维忠实不可约表示 \(\varrho\) 。G 的复化李代数同构于 \(\mathfrak{e}_6\) ，当且仅当表示 \(\operatorname{End}(\varrho)\) 中存在一个 78 维的不可约子表示。
\item
  设 \(q\) 为奇素数的幂，F 为含 \(q\) 个元素的有限域。考虑由 F 上系数的 3 阶上三角幺幂矩阵组成的海森伯群 H \(_3\) (F)。对 \((a, b, c) \in \mathrm{F}^3\) ，我们记
\end{enumerate}

\[
\sigma (a, b, c) = \left( \begin{array}{c c c} 1 & a & c \\ 0 & 1 & b \\ 0 & 0 & 1 \end{array} \right) \in \mathrm{H} _ {3} (\mathrm{F}).
\]

因此映射 \(\sigma\) 是双射。

\begin{enumerate}
\def\labelenumi{\alph{enumi})}
\item
  \(\mathrm{H}_{3}(\mathrm{F})\) 的中心 C 通过把 \(\sigma(a,b,c)\) 映到 c 的群同态与 F 等同。
\item
  存在 \(q^{2}\) 个 \(\mathrm{H}_{3}(\mathrm{F})\) 的一维酉特征标；这些特征标在 C 上是平凡的。
\item
  设 X 为 F 中由元素 \(\sigma(0,b,c)\) （\((b,c)\in\mathrm{F}^{2}\) ）组成的指标为 q 的子群。子群 X 与 \(F^{2}\) 等同。设 \(\psi_{1}\) 与 \(\psi_{2}\) 为 F 的酉特征标，令 \(\psi(\sigma(0,b,c))=\psi_{1}(b)\psi_{2}(c)\) 。酉表示 \(\varrho\) （\(\mathrm{H}_{3}(\mathrm{F})\) 的、由 X 的特征标 \(\psi\) 诱导）是不可约的，当且仅当 \(\psi_{2}\) 非平凡；它在同构意义下只依赖于 \(\psi_{2}\) ，其中心特征标等于 \(\psi_{2}\) 。\(\varrho\) 的特征标的支集等于 C。
\item
  在前面两个问题中构造的不可约表示就是 \(\mathrm{H}_{3}(\mathrm{F})\) 的全部不可约酉表示。
\end{enumerate}

\begin{enumerate}
\def\labelenumi{\alph{enumi})}
\setcounter{enumi}{4}
\tightlist
\item
  设 \(\varrho\) 为 \(\mathrm{H}_{3}(\mathrm{F})\) 的 q 维不可约酉表示之一。表示 \(\varrho\) 是忠实的，且满足 \(\mathrm{M}_{4}(\varrho)=q^{2}\) 。
\end{enumerate}

\begin{enumerate}
\def\labelenumi{\arabic{enumi})}
\setcounter{enumi}{15}
\tightlist
\item
  设 \(r \geqslant 1\) 为整数，G 为 \(\mathbf{U}(\mathbf{C}^{r})\) 的连通紧子群，它在 \(C^{r}\) 上不可约地作用。我们用 D 表示 \(\mathbf{U}(\mathbf{C}^{r})\) 中由对角矩阵组成的子群。
\end{enumerate}

\begin{enumerate}
\def\labelenumi{\alph{enumi})}
\tightlist
\item
  设 A 为 D 的一个含于 G 在 \(\mathbf{U}(\mathbf{C}^{r})\) 中的正规化子内的子群。设 S 为 \(\widehat{A}\) 中由对角系数 \(a = (a_{i,j}) \mapsto a_{i,i}\) （\(1 \leqslant i \leqslant r\) ）组成的子集。假设 S 含有 r 个元素，且方程
\end{enumerate}

\[
\chi_ {1} \chi_ {2} ^ {- 1} = \chi_ {3} \chi_ {4} ^ {- 1}
\]

（其中 \(\chi_{i} \in S\) ）只有形如 \((\chi_{1}, \chi_{1}, \chi_{3}, \chi_{3})\) 与 \((\chi_{1}, \chi_{2}, \chi_{1}, \chi_{2})\) 的解（这时称 \(S\) 是一个 Sidon 集）。

若 G 是半单的，证明 \(\mathbf{G} = \mathbf{S}\mathbf{U}(\mathbf{C}^{r})\) 。（考虑 A 在 \(\operatorname{End}(\mathbf{C}^{r})\) 上的表示，并由假设推出 G 的复化李代数在 D 的共轭下稳定。）

\begin{enumerate}
\def\labelenumi{\alph{enumi})}
\setcounter{enumi}{1}
\tightlist
\item
  群 G 含有 \(\mathbf{SU}(\mathbf{C}^{r})\) ，当且仅当存在元素 \(g \in G\) ，使得 g 的特征值互不相同并在 \(C^{*}\) 中构成一个 Sidon 集。
\end{enumerate}

\begin{enumerate}
\def\labelenumi{\arabic{enumi})}
\setcounter{enumi}{16}
\tightlist
\item
  设 G 为紧群，H 为 G 的闭子群。设 \(\varrho\) 为 H 在有限维希尔伯特空间 E 中的酉表示。
\end{enumerate}

\begin{enumerate}
\def\labelenumi{\alph{enumi})}
\tightlist
\item
  设 \(\pi \in \widehat{\mathbf{G}}\) 。空间 \(\mathrm{Hom}_{\mathrm{G}}(\pi, \mathrm{Ind}_{\mathrm{H}}^{\mathrm{G}}(\varrho))\) 是有限维的。

\end{enumerate}

\begin{enumerate}
\def\labelenumi{\alph{enumi})}
\setcounter{enumi}{1}
\item
  对每个 \(u \in \mathrm{Hom}_{\mathrm{G}}(\pi, \mathrm{Ind}_{\mathrm{H}}^{\mathrm{G}}(\varrho))\) ，\(u\) 的像含于 \(\mathcal{C}(\mathrm{G};\mathrm{E})\) 在 \(\mathrm{L}^2 (\mathrm{G};\mathrm{E})\) 中的像。
\item
  存在线性映射 \(\Phi\) ，它从空间 \(\mathrm{Hom}_{\mathrm{G}}(\pi, \mathrm{Ind}_{\mathrm{H}}^{\mathrm{G}}(\varrho))\) 到空间 \(\mathrm{Hom}_{\mathrm{H}}(\mathrm{Res}_{\mathrm{H}}^{\mathrm{G}}(\pi), \varrho)\) ，使得对每个 \(u \in \mathrm{Hom}_{\mathrm{G}}(\pi, \mathrm{Ind}_{\mathrm{H}}^{\mathrm{G}}(\varrho))\) 与每个 \(x \in \mathrm{E}_{\pi}\) ，\(\Phi(u)(x)\) 是连续函数在 \(e\) 处的值，该连续函数表示 \(u(x)\) 。
\item
  线性映射 \(\Phi\) 是同构。
\item
  假设 \(\varrho\) 不可约。\(\pi\) 在 \(\mathrm{Ind}_{\mathrm{H}}^{\mathrm{G}}(\varrho)\) 中的重数等于 \(\varrho\) 在 \(\mathrm{Res}_{\mathrm{H}}^{\mathrm{G}}(\pi)\) 中的重数（弗罗贝尼乌斯互反公式，参见 A, VIII, p. 392）。
\end{enumerate}

\begin{enumerate}
\def\labelenumi{\arabic{enumi})}
\setcounter{enumi}{17}
\item
  假设 G 是紧的。设 \(\varrho\) 与 \(\varrho'\) 为 G 的酉表示。表示 \(\varrho\) 弱包含于 \(\varrho'\) （参见 V, p. 497 的习题 11），当且仅当对 G 的每个不可约酉表示 \(\pi\) ，条件 \(m_{\pi}(\varrho) \geqslant 1\) 蕴含 \(m_{\pi}(\varrho') \geqslant 1\) 。（化归到 \(\varrho'\) 不可约的情形；对每个规范化的对角矩阵系数 \(f\) （属于 \(\varrho\) ），应用 V, p.~498 的习题 13 证明存在 \(\varrho'\) 的一个与 \(f\) 不正交的对角矩阵系数，再用正交关系得出结论。）
\item
  设 G 为赋有规范化哈尔测度 \(\mu\) 的紧群，\(f \in \mathcal{C}(\mathrm{G})\) 为 G 上的连续函数。存在不可约表示 \(\pi \in \widehat{G}\) ，使得 \(f = \dim(\pi)^{-1}\chi_{\pi}\) ，当且仅当
\end{enumerate}

\[
\int_ {\mathrm{G}} f (x g x ^ {- 1} h) d \mu (x) = f (g) f (h)
\]

对所有 \((g,h)\in \mathrm{G}^2\) 成立

\begin{enumerate}
\def\labelenumi{\arabic{enumi})}
\setcounter{enumi}{19}
\tightlist
\item
  设 H 为赋有左哈尔测度 \(\nu\) 的局部紧拓扑群。设 G 为 H 的紧子群。设 \(\varrho\) 为 H 在希尔伯特空间 E 中的不可约酉表示。
\end{enumerate}

我们假设 \(\pi\) 在 \(\varrho\) 到 G 的限制中的重数对每个不可约表示 \(\pi \in \widehat{G}\) 都是有限的。我们把 G 上的函数与 H 上紧支集的测度等同起来。

\begin{enumerate}
\def\labelenumi{\alph{enumi})}
\item
  对每个 \(\pi \in \widehat{\mathrm{G}}\) ，E 的自同态 \(\varrho(f * \chi_{\pi})\) 是有限秩的。
\item
  对每个 \(f \in \mathrm{L}^{1}(\mathrm{H})\) ，自同态 \(\varrho(f)\) 是紧的。
\end{enumerate}

\begin{enumerate}
\def\labelenumi{\arabic{enumi})}
\setcounter{enumi}{20}
\tightlist
\item
  设 \(n \geqslant 1\) 为整数。
\end{enumerate}

\begin{enumerate}
\def\labelenumi{\alph{enumi})}
\tightlist
\item
  对每个整数 \(r > n^{2}\) 与每个族 \((a_{1}, \ldots, a_{r})\) （\(\mathcal{L}(\mathbf{C}^{n})\) 的元素），我们有
\end{enumerate}

\[
\sum_ {\sigma \in \mathfrak {S} _ {r}} \varepsilon (\sigma) a _ {\sigma (1)} \dots a _ {\sigma (r)} = 0.\tag{4}
\]

\begin{enumerate}
\def\labelenumi{\alph{enumi})}
\setcounter{enumi}{1}
\tightlist
\item
  我们用 \(r(n)\) 表示最小的整数 \(r \geqslant 1\) ，使公式 (4) 对每个族 \((a_i)_{1 \leqslant i \leqslant r}\) 成立。我们有 \(r(n) \geqslant r(n-1) + 2\) 。（考虑一族 \((a_1, \ldots, a_{r(n-1)-1})\) （属于 \(\mathcal{L}(\mathbf{C}^{n-1})\) ），使得记作 b 的上述表达式非零；把 \(\mathcal{L}(\mathbf{C}^n)\) 的元素解释为矩阵，选取 i 与 j 使 \(b_{i,j}\) 非零；把诸 \(a_i\) 与 b 视为 \(\mathcal{L}(\mathbf{C}^n)\) 的元素（通过添加一行一列零），然后定义矩阵 \(a_{r(n-1)}\) 与 \(a_{r(n-1)+1}\) ，使得 \((a_1, \ldots, a_{r(n-1)+1})\) 的表达式 (4) 非零。）
\end{enumerate}

\begin{enumerate}
\def\labelenumi{\arabic{enumi})}
\setcounter{enumi}{21}
\tightlist
\item
  设 \(r \geqslant 1\) 为整数，\(\mathrm{H} \subset \mathbf{GL}(r, \mathbf{C})\) 为闭子群。
\end{enumerate}

\begin{enumerate}
\def\labelenumi{\alph{enumi})}
\item
  H 上形如 \(g \mapsto \langle \varrho(g)x, \lambda \rangle\) 的函数所成的集合（其中 \(\varrho\) 是 H 在有限维复向量空间 E 中的连续线性表示，\(x \in E\) ，\(\lambda \in E'\) ）是 \(\mathcal{C}(\mathrm{H})\) 的稠密对合子代数。
\item
  对任意非零的函数 \(f \in \mathcal{K}(\mathrm{H})\) ，存在 H 在有限维复向量空间中的一个连续线性表示 \(\varrho\) ，使得 \(\varrho(f)\) 非零。
\item
  假设 H 是连通半单实李群。对每个非零的 \(f \in \mathcal{K}(H)\) ，存在 H 在有限维复向量空间中的不可约连续线性表示 \(\varrho\) ，使得 \(\varrho(f)\) 非零。
\end{enumerate}

\begin{enumerate}
\def\labelenumi{\arabic{enumi})}
\setcounter{enumi}{22}
\tightlist
\item
  我们沿用上一习题的假设与记号，并假设 H 是连通半单实李群。
\end{enumerate}

设 G 为 H 的紧子群。假设存在 H 的连通可解闭子群 N，使得 H = GN。（对每个半单实李群，当 G 是 H 的极大紧子群时这一条件满足；例如当 G = SL(n, R) 时，参见 INT, VII, p.~91, § 3, n° 3, 命题 7。）

设 \(\pi \in \widehat{\mathbf{G}}\) ，\(e_{\pi} = \dim (\pi)\overline{\chi}_{\pi}\) 。设 \(\varrho\) 为上一习题最后一问给出的 H 的不可约表示；我们用 E 表示 \(\varrho\) 的空间。

\begin{enumerate}
\def\labelenumi{\alph{enumi})}
\item
  \(\varrho\) 到 G 的限制存在循环向量。（利用李定理，LIE, III, p.~241, 推论，找出非零的 \(x \in E\) ，使 C \(x\) 在 N 的作用下不变。）\(\pi\) 在 \(\varrho\) 到 G 的限制中的重数至多等于 \(\dim(\pi)\) 。
\item
  设 \(r\) 为习题 22 中定义的整数 \(r(\dim (\pi)^2)\) ；恒等式
\end{enumerate}

\[
\sum_ {\sigma \in \mathfrak {S} _ {r}} \varepsilon (\sigma) a _ {\sigma (1)} \dots a _ {\sigma (r)} = 0
\]

对每个元素族 \(a_{i}=e_{\pi}*g_{i}*e_{\pi}\) （属于 H 上紧支集测度的代数 \(\mathcal{M}_{c}(\mathrm{H})\) ，其中 \(g_{i}\in\mathcal{K}(\mathrm{H})\) （对 \(1\leqslant i\leqslant r\) ））成立。

\begin{enumerate}
\def\labelenumi{\alph{enumi})}
\setcounter{enumi}{2}
\tightlist
\item
  设 \(\varpi\) 为 H 在希尔伯特空间 F 中的不可约酉表示。设 \(F_{\pi}\) 为 \(\pi\)-同型分量（在 \(\varpi\) 到 G 的限制中），\(p\) 为 F 的以 \(F_{\pi}\) 为像的正交投影算子。对每个族 \((u_i)_{1 \leqslant i \leqslant r}\) （属于 \(\mathcal{L}(F_{\pi})\) ），我们有
\end{enumerate}

\[
\sum_ {\sigma \in \mathfrak {S} _ {r}} \varepsilon (\sigma) u _ {\sigma (1)} \dots u _ {\sigma (r)} = 0
\]

  （利用 V, p.~486 的习题 1）。

\begin{enumerate}
\def\labelenumi{\alph{enumi})}
\setcounter{enumi}{3}
\tightlist
\item
  \(\pi\) 在 \(\varpi\) 到 G 的限制中的重数至多等于 \(\dim(\pi)\) 。（利用 V, p.~510 的习题 22 的结果。）
\end{enumerate}

\begin{enumerate}
\def\labelenumi{\arabic{enumi})}
\setcounter{enumi}{23}
\tightlist
\item
  称拓扑群 G 是极小概周期的，如果从 G 到紧群的每个连续同态都是平凡的。
\end{enumerate}

\begin{enumerate}
\def\labelenumi{\alph{enumi})}
\item
  拓扑群 G 是极小概周期的，当且仅当 G 不容许非平凡的有限维酉表示。若 G 是交换的，则 G 是极小概周期的，当且仅当 G 的每个连续酉特征标都是平凡的。
\item
  设 G 为拓扑群，\(g \in G\) 满足：对每个整数 \(k \geqslant 1\) ，存在整数 \(m \geqslant 1\) ，使得 \(g^{km}\) 共轭于 g。则有 \(\varrho(g) = 1_{\mathrm{E}}\) 对 G 在希尔伯特空间 E 中的每个有限维酉表示 \(\varrho\) 成立。
\item
  设 k 为 C 的子域。赋有离散拓扑的群 \(\mathbf{SL}(2,k)\) 是极小概周期的。（验证上一问题适用于 \(\begin{pmatrix}1 & z \\ 0 & 1\end{pmatrix}\) （对所有 \(z \in k\) ）。）
\end{enumerate}

\begin{enumerate}
\def\labelenumi{\arabic{enumi})}
\setcounter{enumi}{24}
\tightlist
\item
  设 G 为赋有左哈尔测度 \(\nu\) 的局部紧群，\(K \subset G\) 为赋有规范化哈尔测度 \(\mu\) 的 G 的紧子群。
\end{enumerate}

我们用 \(\mathcal{H}(\mathrm{G},\mathrm{K})\) 表示由函数 \(\varphi \in \mathcal{H}(\mathrm{G})\) （满足 \(\varphi (k_1gk_2) = \varphi (g)\) ，对所有 \(g\in \mathrm{G}\) 与每个 \((k_{1},k_{2})\in \mathrm{K}\times \mathrm{K}\) ）所成的空间。

空间 \(\mathcal{H}(\mathrm{G},\mathrm{K})\) 是代数 \(\mathcal{M}_c(\mathrm{G})\) 的一个含幺子代数；我们称 \((\mathrm{G},\mathrm{K})\) 为盖尔范德对，如果代数 \(\mathcal{H}(\mathrm{G},\mathrm{K})\) 是交换的。\(a)\) 映射 \(u^{\mathrm{K}}\) 把 \(f\) 对应到函数 \(f^{\mathrm{K}}\) （由

\[
f ^ {\mathrm{K}} (g) = \int_ {\mathrm{K} \times \mathrm{K}} f (k _ {1} g k _ {2}) d (\mu \otimes \mu) (k _ {1}, k _ {2})
\]

所定义）；它是 \(\mathcal{K}(\mathrm{G})\) 中的一个投影算子，其像为 \(\mathcal{H}(\mathrm{G},\mathrm{K})\) 。

\begin{enumerate}
\def\labelenumi{\alph{enumi})}
\setcounter{enumi}{1}
\tightlist
\item
  设 \(f\in \mathcal{K}(\mathrm{G})\) 。我们有
\end{enumerate}

\[
\int_ {\mathrm{G}} f (g) d \nu (g) = \int_ {\mathrm{G}} f ^ {\mathrm{K}} (g) d \nu (g)
\]

并且 \((f\varphi)^{\mathrm{K}} = f^{\mathrm{K}}\varphi\) 对每个函数 \(\varphi \in \mathcal{H}(\mathrm{G},\mathrm{K})\) 成立。

\begin{enumerate}
\def\labelenumi{\alph{enumi})}
\setcounter{enumi}{2}
\item
  设 \(C \subset G\) 为紧子集。存在函数 \(\varphi \in \mathcal{H}(G, K)\) ，使得 \(\varphi(x) = 1\) 对所有 \(x \in C\) 成立。
\item
  设 (G,K) 为盖尔范德对。群 G 是幺模的。（注意 G 的模属于 \(\mathcal{H}(\mathrm{G},\mathrm{K})\) ，并把 \(\varphi \in \mathcal{H}(\mathrm{G},\mathrm{K})\) 的积分解释为适当选取的卷积在 \(e\) 处的值。）
\item
  假设存在连续对合 \(\theta\colon\mathrm{G}\to\mathrm{G}\) ，使得 \(\mathrm{K}\theta(x)\mathrm{K}=\mathrm{K}x^{-1}\mathrm{K}\) 对所有 \(x\in\mathrm{G}\) 成立。则 (G,K) 是盖尔范德对。（验证映射 \(f\mapsto\nu(f\circ\theta)\) （在 \(\mathcal{K}(\mathrm{G})\) 上）是 G 上的一个左哈尔测度，进而它等于 \(\underline{\nu}\) ；推出 \((f*g)\circ\theta=(f\circ\theta)*(g\circ\theta)\) ，并与公式 \(f*g=\check{g}*f\) 比较。）
\item
  若 G 是紧的，则 \((\mathbf{G} \times \mathbf{G}, \Delta)\) 是盖尔范德对，其中 \(\Delta\) 是 \(G \times G\) 的对角子群，即同态 \(x \mapsto (x, x)\) （从 G 到 \(G \times G\) ）的像。
\end{enumerate}

\begin{enumerate}
\def\labelenumi{\arabic{enumi})}
\setcounter{enumi}{25}
\tightlist
\item
  我们继续沿用上一习题的记号。我们用 \(\varphi_{\mathrm{K}}\) 表示 K 的特征函数。
\end{enumerate}

\begin{enumerate}
\def\labelenumi{\alph{enumi})}
\item
  设 \(\pi\) 为 G 在复希尔伯特空间 E 中的酉表示。\(\mathcal{M}_c(\mathrm{G})\) 在 E 中的表示通过过渡到子空间，定义了代数 \(\mathcal{H}(\mathrm{G},\mathrm{K})\) 在 E 的 K 不变向量所成空间 \(\mathrm{E}^{\mathrm{K}}\) 中的表示。
\item
  若 E 容许循环向量 x，使得 \(x \in E^{K}\) ，则 \(\pi\) 是不可约的。
\item
  假设 \(\pi\) 不可约。若空间 \(E^{K}\) 不约化为 0，则 \(\mathcal{H}(G,K)\) 在 \(E^{K}\) 中的表示是不可约的。
\end{enumerate}

\begin{enumerate}
\def\labelenumi{\alph{enumi})}
\setcounter{enumi}{3}
\tightlist
\item
  设 \(\varphi\) 为 G 上的正定函数，\((\varrho, x)\) 为 \(\varphi\) 的一个循环希尔伯特实现。我们有 \(\varphi \in \mathcal{H}(\mathrm{G}, \mathrm{K})\) ，当且仅当 \(x \in \mathrm{E}^{\mathrm{K}}\) 。
\end{enumerate}

\begin{enumerate}
\def\labelenumi{\arabic{enumi})}
\setcounter{enumi}{26}
\tightlist
\item
  我们继续沿用上一习题的记号。假设 (G, K) 是盖尔范德对。
\end{enumerate}

我们称 \(\varphi\in\mathcal{H}(\mathrm{G},\mathrm{K})\) 是 K-球函数，如果映射

\[
f \mapsto \int_ {\mathrm{G}} f \check {\varphi} d \nu
\]

是代数 \(\mathcal{H}(\mathrm{G},\mathrm{K})\) 的一个非零特征标。

\begin{enumerate}
\def\labelenumi{\alph{enumi})}
\tightlist
\item
  设 \(\varphi \in \mathcal{H}(\mathrm{G},\mathrm{K})\) 。证明下列条件等价：
\end{enumerate}

\begin{enumerate}
\def\labelenumi{(\roman{enumi})}
\item
  函数 \(\varphi\) 是 K-球函数；
\item
  对所有 \(x\) 与 \(y\) （属于 \(\mathrm{G}\) ），我们有
\end{enumerate}

\[
\int_ {\mathrm{K}} \varphi (x k y) d \mu (k) = \varphi (x) \varphi (y);
\]

\begin{enumerate}
\def\labelenumi{(\roman{enumi})}
\setcounter{enumi}{2}
\tightlist
\item
  有 \(\varphi(e)=1\) ，且对每个函数 \(\psi\in\mathcal{H}(\mathrm{G},\mathrm{K})\) ，存在 \(c\in\mathbf{C}\) 使得 \(\psi*\varphi=c\varphi\) 。
\end{enumerate}

\begin{enumerate}
\def\labelenumi{\alph{enumi})}
\setcounter{enumi}{1}
\item
  我们用 \(\mathcal{H}^1 (\mathrm{G},\mathrm{K})\) 表示由 \(f\in \mathrm{L}^1 (\mathrm{G})\) （满足 \(\varrho_{\mathrm{G}}(k_1,k_2)f = f\) ，对所有 \((k_{1},k_{2})\in \mathrm{K}\times \mathrm{K}\) ）所成的空间。它是对合巴拿赫代数 \(\mathrm{L}^1 (\mathrm{G})\) 的一个巴拿赫子代数；它是交换的。
\item
  代数 \(\mathcal{H}(\mathrm{G},\mathrm{K})\) 与 \(\mathcal{H}^1 (\mathrm{G},\mathrm{K})\) 的一个稠密子代数等同。\\
\item
  把 \(\varphi \in \mathcal{H}(\mathrm{G},\mathrm{K})\) 对应到线性型
\end{enumerate}

\[
f \mapsto \int_ {\mathrm{G}} f \check {\varphi} d \nu
\]

的映射（在 \(\mathcal{H}^{1}(\mathrm{G},\mathrm{K})\) 上）定义了从 G 上有界 K-球函数所成的空间到 \(\mathcal{H}^{1}(\mathrm{G},\mathrm{K})\) 的非零特征标所成空间上的一个双射。（注意 \(\mathcal{H}^{1}(\mathrm{G},\mathrm{K})\) 的每个特征标都是连续的且范数 \(\leqslant 1\) ，参见 I, p.~104 的命题 2。）

\begin{enumerate}
\def\labelenumi{\arabic{enumi})}
\setcounter{enumi}{27}
\tightlist
\item
  设 G 为局部紧群，K 为 G 的紧子群。
\end{enumerate}

\begin{enumerate}
\def\labelenumi{\alph{enumi})}
\item
  偶 (G, K) 是盖尔范德对，当且仅当对 G 在复希尔伯特空间 E 中的每个酉表示 \(\pi\) ，\(E^{K}\) 的维数至多为 1。（为证该条件是必要的，证明交换对合巴拿赫代数在希尔伯特空间中的不可约表示是一维的；为证它是充分的，利用 V, p.~454 的定理 4 的推论验证代数 \(\mathcal{H}(G,K)\) 的特征标分离其点。）
\item
  假设 (G, K) 是盖尔范德对。设 \(\varphi \in \mathcal{H}(\mathrm{G},\mathrm{K})\) 。假设 \(\varphi\) 是正定函数且 \(\varphi(e) = 1\) 。证明 \(\varphi\) 是 K-球函数，当且仅当 \(\varphi\) 有一个希尔伯特表示 \((\varrho, x)\) ，使得 \(\varrho\) 不可约。
\item
  假设 \((\mathrm{G},\mathrm{K})\) 是盖尔范德对。G 上正定 K-球函数所成的集合与由不可约表示 \(\pi\in\widehat{G}\) （满足 \(E_{\pi}^{K}\) 非零）所成的集合双射对应。
\end{enumerate}

\begin{enumerate}
\def\labelenumi{\arabic{enumi})}
\setcounter{enumi}{28}
\tightlist
\item
  设 G 为紧群，K 为 G 的闭子群。假设 (G, K) 是盖尔范德对。我们用 \(\nu\) （相应地 \(\mu\) ）表示 G（相应地 K）上的规范化哈尔测度。
\end{enumerate}

设 \(\varphi\) 为 G 上的一个 K-球函数。

\begin{enumerate}
\def\labelenumi{\alph{enumi})}
\item
  我们用 \(E_{\varphi}\) 表示 \(\mathcal{C}(G)\) 中形如 \(\alpha * \varphi\) 的函数所成的子空间，其中 \(\alpha\) 是 G 上有限支集的测度。对每个 \(\psi \in \mathcal{H}(G, K)\) ，自同态 \(v_{\psi}: f \mapsto f * \psi\) （\(\mathcal{C}(G)\) 的）是紧的，且 \(E_{\varphi}\) 含于 \(v_{\psi}\) 的一个特征空间。
\item
  空间 \(E_{\varphi}\) 是有限维的。它与正则表示 \(\gamma_{G}\) 的 G-有限向量所成空间的一个含 \(\varphi\) 的子空间等同。
\item
  \(E_{\varphi}\) 的 K 不变向量所成的空间是一维的，且 G 在 \(E_{\varphi}\) 中的表示是不可约的。
\end{enumerate}

\begin{enumerate}
\def\labelenumi{\alph{enumi})}
\setcounter{enumi}{3}
\tightlist
\item
  设 \((f_i)_{i\in I}\) 为 \(\mathrm{E}_{\varphi}\) 的规范正交基，\(k\colon \mathbf{G}\times \mathbf{G}\to \mathbf{C}\) 为由

\[
k (x, y) = \sum_ {i \in \mathrm{I}} \overline {{f _ {i} (x)}} f _ {i} (y).
\]

定义的函数。有

\[
f (x) = \int_ {\mathrm{G}} k (x, y) f (y) d \nu (y)
\]

对所有 \(f \in \mathrm{E}_{\varphi}\) 与 \(x \in \mathrm{G}\) 成立。

\end{enumerate}

\begin{enumerate}
\def\labelenumi{\alph{enumi})}
\setcounter{enumi}{4}
\item
  有 \(k(x,y)=\dim(\mathrm{E}_{\varphi})\varphi(y^{-1}x)\) 对所有 \((x,y)\in\mathrm{G}\times\mathrm{G}\) 成立。
\item
  函数 \(\varphi\) 是正定型的。
\item
  设 \(\varphi\) 为 G 上的 K-球函数，\((\pi, x)\) 为 \(\varphi\) 的一个希尔伯特实现。表示 \(\pi\) 是不可约的；我们有
\end{enumerate}

\[
\varphi (x) = \int_ {\mathrm{K}} \chi_ {\pi} (x ^ {- 1} k) d \mu (k)
\]

对所有 \(x \in \mathrm{G}\) 成立，且 \(\int_{\mathrm{G}} |\varphi|^2 d\nu = \dim(\pi)^{-1}\) 。

\begin{enumerate}
\def\labelenumi{\arabic{enumi})}
\setcounter{enumi}{29}
\tightlist
\item
  设 \(n \in N\) 。用 G 表示仿射群 \(\mathrm{A}(\mathbf{R}^{n})\) （空间 \(R^{n}\) 的）中由元素 \(g \in \mathrm{A}(\mathbf{R}^{n})\) （其相伴线性映射 \(g_{0}\) 属于 \(\mathbf{SO}(\mathbf{R}^{n})\) ）所成的子群；它是一个实李群（参见 A, II, p. 131 与 LIE, III, p. 102, 例）。我们把 G 的元素记作 \(g = (g_{0}, a)\) ，其中 \(g_{0} \in \mathbf{SO}(\mathbf{R}^{n})\) ，\(a \in R^{n}\) 。
\end{enumerate}

\begin{enumerate}
\def\labelenumi{\alph{enumi})}
\item
  偶 \((\mathbf{G},\mathbf{SO}(\mathbf{R}^n))\) 是盖尔范德对。
\item
  设 \(\varphi \in \mathcal{H}(\mathrm{G},\mathbf{SO}(\mathbf{R}^n))\) 。用 \(f\colon \mathbf{R}^n\to \mathbf{C}\) 表示由 \(f(r) = \varphi (1_{\mathbf{R}^N},g(e_1))\) 定义的函数。映射 \(u\colon \varphi \mapsto f\) 是从 \(\mathcal{H}(\mathrm{G},\mathrm{K})\) 到 \(\mathcal{K}(\mathbf{R}^n)\) 的一个单射复代数态射；其像是 \(\mathcal{K}(\mathbf{R}^n)\) 中由函数 \(f\) （满足 \(f(g_0(x)) = f(x)\) ，对所有 \(g_0\in \mathbf{SO}(\mathbf{R}^n)\) 与每个 \(x\in \mathbf{R}^n\) ）组成的子空间。
\item
  我们用 \(\Delta\) 表示 \(\mathbf{R}^n\) 上的拉普拉斯算子。函数 \(\varphi \in \mathcal{K}(\mathrm{G})\) 是球函数，当且仅当函数 \(f = u(\varphi)\) 是 \(\mathrm{C}^{\infty}\) 类的、满足 \(f(0) = 1\) ，且存在 \(\lambda \in \mathbf{C}\) 使得 \(\Delta f = \lambda f\) 。（利用 \(\Delta(f \circ g_0) = (\Delta f) \circ g_0\) 对所有 \(g_0 \in \mathbf{SO}(\mathbf{R}^n)\) 成立这一事实。）
\end{enumerate}

\begin{enumerate}
\def\labelenumi{\arabic{enumi})}
\setcounter{enumi}{30}
\tightlist
\item
  设 \(n\) 为整数 \(\geqslant 1\) 。我们用 N 表示 \(\mathbf{R}^n\) 上的欧几里得范数，用 \(\lambda_n\) 表示 \(\mathbf{R}^n\) 上的勒贝格测度。我们用 \(\omega_n\) 表示球面 \(\mathbf{S}_n \subset \mathbf{R}^{n+1}\) 上唯一的、在正交群 \(\mathbf{O}(n+1, \mathbf{R})\) 的作用下不变的、总质量为 1 的正测度（参见 INT, VIII, p.~116, 习题 8）。
\end{enumerate}

我们把 \(R^{n}\) 与射影空间 \(\mathbf{P}_{n}(\mathbf{R})\) 中射影超平面 H 的余集等同起来（参见 TG, VI, p.~15, 命题 5），并把 \(R^{n}\) 上的测度解释为 \(\mathbf{P}_{n}(\mathbf{R})\) 上支集含于 H 的余集中的测度。

对每个 \(g \in \mathbf{GL}(n+1, \mathbf{R})\) ，我们用 \(\widetilde{g}\) 表示由 g 定义的 \(\mathbf{P}_{n}(\mathbf{R})\) 的同胚。

\begin{enumerate}
\def\labelenumi{\alph{enumi})}
\tightlist
\item
  正测度 \(c_{n}\) （在 \(\mathbf{R}^n\) 上）由
\end{enumerate}

\[
c _ {n} = \pi^ {- (n + 1) / 2} \Gamma \Big (\frac {n + 1}{2} \Big) (1 + \mathrm{N} ^ {2}) ^ {- (n + 1) / 2} \lambda_ {n}
\]

定义的（“柯西测度”）是有界的且总质量为 1。若 \(n = 1\) ，它就是参数 1 的稳定测度（参见习题 V, p.~496）。对每个射影超平面 \(\mathrm{H}'\) （\(\mathbf{P}_n(\mathbf{R})\) 的），我们有 \(c_n(\mathrm{H}') = 0\) 。

\begin{enumerate}
\def\labelenumi{\alph{enumi})}
\setcounter{enumi}{1}
\item
  从 \(\mathbf{R}^{n+1} - \{0\}\) 到 \(\mathbf{P}_n(\mathbf{R})\) 的典范映射通过过渡到子空间，定义了连续满射 \(\sigma: \mathbf{S}_n - (\mathbf{R}^n \times \{0\}) \to \mathbf{R}^n\) 。我们有 \(\sigma(\omega_n) = c_n\) 。
\item
  设 \(g \in \mathbf{GL}(n+1, \mathbf{R})\) 。存在 \(a \in \mathbf{GL}(n, \mathbf{R})\) 与 \(b \in R^n\) ，使得像测度 \(\widetilde{g}(c_n)\) 等于 \(c_n\) 在仿射映射 \(x \mapsto a(x) + b\) 下的像。
\end{enumerate}

在下面，我们打算证明这一断言的一个逆命题。

\begin{enumerate}
\def\labelenumi{\alph{enumi})}
\setcounter{enumi}{3}
\tightlist
\item
  设 \(\mu\) 为 \(\mathbf{P}_{n}(\mathbf{R})\) 上的有界正测度，使得每个射影超平面都是 \(\mu\) -可忽略的。由 \(g \in \mathbf{GL}(n+1, \mathbf{R})\) （满足 \(g(\mu) = \mu\) 且 \(|\det(g)| = 1\) ）所成的集合是 \(\mathbf{GL}(n+1, \mathbf{R})\) 的一个紧子群。

\end{enumerate}

\begin{enumerate}
\def\labelenumi{\alph{enumi})}
\setcounter{enumi}{4}
\item
  设 \(\mu\) 为 \(\mathbf{P}_n(\mathbf{R})\) 上总质量为 1 的有界正测度，使得对每个 \(g \in \mathbf{GL}(n + 1, \mathbf{R})\) ，存在仿射映射 \(f\) （\(\mathbf{R}^n\) 的）满足 \(g(\mu) = f(\mu)\) 。每个射影超平面都是 \(\mu\) -可忽略的。
\item
  设 \(K_{\mu}\) 为由 \(g \in \mathbf{GL}(n+1, \mathbf{R})\) （满足 \(g(\mu) = \mu\) ）所成的子群。存在 \(h \in \mathbf{GL}(n+1, \mathbf{R})\) ，使得 \(K_{\mu} \subset h \mathbf{O}(n+1, \mathbf{R}) h^{-1}\) 。
\item
  设 \(K = h^{-1}K_{\mu}h\) 。我们有 \(\mathbf{O}(n + 1, \mathbf{R}) = \mathbf{G} \cdot \mathbf{K}\) ，其中 G 是 \(\mathbf{O}(n + 1, \mathbf{R})\) 中由矩阵
\end{enumerate}

\[
\left( \begin{array}{c c} x & 0 \\ 0 & 1 \end{array} \right)
\]

（\(x\in \mathbf{O}(n,\mathbf{R})\) ）组成的子群

\begin{enumerate}
\def\labelenumi{\alph{enumi})}
\setcounter{enumi}{7}
\tightlist
\item
  测度 \(h^{-1}(\mu)\) 与 \(c_{n}\) 重合。（注意群 K 传递地作用于 \(S_{n}\) ，且像测度 \(\sigma(h^{-1}(\mu))\) 是 \(S_{n}\) 上的一个 K 不变测度。）
\end{enumerate}

\begin{enumerate}
\def\labelenumi{\arabic{enumi})}
\setcounter{enumi}{31}
\tightlist
\item
  设 G 为局部紧群。假设存在 G 的一个既紧又开的子群 H。我们用 \(\mu\) 表示 G 上满足 \(\mu(\mathrm{H}) = 1\) 的唯一的左哈尔测度，用 \(\varphi\) 表示 H 的特征函数。
\end{enumerate}

\begin{enumerate}
\def\labelenumi{\alph{enumi})}
\item
  设 \(\pi\) 为 G 在复希尔伯特空间 E 中的一个平方可积表示。证明自同态 \(\pi(\varphi)\) 是希尔伯特–施密特自同态。
\item
  空间 \(E^{H}\) 是有限维的，且满足
\end{enumerate}

\[
\dim (\mathrm{E} ^ {\mathrm{H}}) \leqslant \frac {1}{c (\pi)}
\]

其中 \(c(\pi)\) 是 \(\pi\) 相对于 \(\{e\}\) 的形式次数。（利用 V, p.~487 的习题 6。）

\begin{enumerate}
\def\labelenumi{\alph{enumi})}
\setcounter{enumi}{2}
\tightlist
\item
  离散无限群没有平方可积表示；若 \(n \geqslant 2\) 为整数，则群 \(\mathbf{SL}(n, \mathbf{Z})\) 没有模中心的平方可积表示。
\end{enumerate}

\specialchapter{历史注记}

（第 I 至 V 章）

特征值与傅里叶级数的概念在 1830 年已为人熟知，并贯穿十九世纪的许多分析领域（参见~ÉHM, pp.~114, 260, 275）。但谱理论与调和分析——按我们在这本书中所理解的意义——直到大约一个世纪之后才开始具有其现在的形式；这与它们在 1930 年前后同群表示研究的汇合、随后在 1940 年前后同赋范代数研究的汇合是同时发生的。

我们在这里仅限于勾画它们从 1906 年（希尔伯特在其积分方程工作中引入“连续谱”概念之时）到 1945–1950 年间演化的主要脉络；在此期间，这里所陈述的理论基本上获得了保持至今的形式。

\histitem{连续谱的发现}

在关于拓扑向量空间的注记中（参见 ÉHM, pp. 262–263），我们描述了 I. 弗雷德霍姆如何在 1900 年前后——继承 C. 诺伊曼、V. 沃尔泰拉与 H. 庞加莱的工作——开始考虑方程

\[
u (x) - \lambda \int_ {\mathrm{I}} \mathrm{K} (x, y) u (y) d y = f (x)\tag{1}
\]

其中未知函数为 \(u\colon I\to\mathbf{C}\) ，这里区间 \(I\subset\mathbf{R}\) 是紧的，且核 \(K\) 在 \(I\times I\) 上连续。他巧妙地使用“无穷行列式”——其基础是庞加莱与 H. 冯·科赫的思想——得到了对复变量 \(\lambda\) 的依赖为亚纯的解族，进而证明了关于解的存在的著名“择一定理”（参见 III, p.~81, 定理 5）。但弗雷德霍姆随即把他的结果特殊化到 \(\lambda = -1\) 的情形，而没有利用方程 (1) 是一个特征值问题 \(^{(1)}\) 这一事实。

当核 K 对称时，我们还看到希尔伯特在其关于积分方程的第一篇论文 [31] 中如何认识到弗雷德霍姆问题与求二次型的“主轴”之间的亲缘关系。从核 K 的有限“截段”（离散化）的这种约化出发，他通过取极限得到公式

\[
\int_ {\mathrm{I}} \mathrm{K} (s, t) x (s) x (t) d t = \sum_ {n = 1} ^ {\infty} \frac {1}{\lambda_ {n}} \int_ {\mathrm{I}} \varphi_ {n} (s) x (s) d s,\tag{2}
\]

其中 \(\lambda_{n}\) 是核 K 的特征值，\(\varphi_{n}\) 构成相伴特征向量 \(^{(2)}\) 的规范正交系，并且只要 x 在 I 上平方可积等式就成立（参见~ÉHM, p.~264）。

希尔伯特面临的主要障碍涉及从有限和过渡到公式 (2) 的积分。他从这一公式推出了特征值的存在性——不过，他受有限维中确定特征值的经典方法的启发（参见~IV, p.~153 第 4 目），为后者给出了一个与取极限无关的变分刻画。E. 施密特在 1905 年将证明，借鉴 H. 施瓦茨的工作，如何不经过弗雷德霍姆与希尔伯特的方法所依赖的“无穷行列式”而得到 (2)（参见~ÉHM, p.~265）。

但在 1906 年，希尔伯特 [32] 已转向约化二次型

\[
\mathrm{B} (x, x) = \sum_ {p, q = 0} ^ {\infty} b _ {p q} x _ {p} x _ {q}, \qquad x = (x _ {p}) _ {p \in \mathbf {N}}
\]

这一更一般的问题，二次型定义在 \(E = \ell_{\mathbf{C}}^{2}(\mathbf{N})\) 上。他心里显然想着：通过过渡到 E 的规范正交基下的系数，问题 (1) 可化为求“无穷线性方程组”

\[
x _ {p} + \sum_ {q = 0} ^ {\infty} b _ {p q} x _ {q} = f _ {p} \quad (p \in \mathbf {N}).
\]

的解。然而他注意到，弗雷德霍姆问题所处理的是非常特殊的二次型：\(\mathrm{B}(x_{n}, x_{n})\) 趋于 \(\mathrm{B}(x, x)\) （只要 \(x_{n}\) 弱收敛于 x）；他称之为“全连续的”（vollstetig）。

希尔伯特强调它们与 E 上有界的二次型之间的本质差别，并决定着手对后者作更一般的约化。他的方法依旧是从“有限截段”

\[
(x _ {0}, \ldots , x _ {n}) \mapsto \sum_ {p, q = 0} ^ {n} b _ {p q} x _ {p} x _ {q}
\]

的约化出发，然后让 n 趋于无穷。为完成这一取极限的过程，他利用了 T. 斯蒂尔杰斯在其关于连分数的工作 [71] 中引入的积分思想。他证明每个有界二次型经过变量的一个正交变换后都可以写成

\[
\sum_ {p \in {\bf N}} \frac {1}{\lambda_ {p}} x _ {p} ^ {2} + \int_ {s \in {\bf R}} \frac {1}{s} d \sigma (s, x).\tag{3}
\]

的形式。在这个公式中，\(\lambda_{p}\) 是 B 的特征值，\(x\mapsto\sigma(s,x)\) （对固定的 \(s\in\mathbf{R}\) ）是 E 上正定的分离二次型，而 \(s\mapsto\sigma(s,x)\) \(^{(3)}\) （对固定的 \(x\in\mathrm{E}\) ）是连续的递增函数，它在 \(-\infty\) 处趋于 0，其极限值 \(\|x\|^2\) 在 \(+\infty\) 处达到（斯蒂尔杰斯记号参见 ÉHM, p. 284）。

用 S 表示 \(\overline{R}\) 中这样的点所成的集合：不存在使 \(s \mapsto \sigma(s, x)\) 对所有 x 都保持常数的邻域；于是 (3) 中的积分自然可以在 S 上取。希尔伯特把集合 S 称为“连续谱”（Streckenspektrum），把 B 的特征值所成的集合称为“点谱”（Punktspektrum），把两者的并称为“谱”（Spektrum）。这个来自光学术语的词，曾在 1897 年被 W. 维尔丁格 [91] 用于带周期系数的微分方程的研究。

希尔伯特很快便从他的“谱”方法中为经典分析的若干问题汲取了深刻的革新，这些问题涉及常微分方程（特别是 Sturm–Liouville 型）、偏微分方程或复变量函数的研究。我们在这里不复述二十世纪头二十五年间随之而来的丰富成果，而请读者参看 E. 海林格与 O. 托普利茨的综述 [29]。不过我们仍要提到其中一些关于积分方程的工作，在它们之中可以看到抽象理论后来发展的萌芽。

对希尔伯特及其学生而言——他们忠实于线性代数中的德国传统——模型始终是主轴定理；因此所研究的总是双线性型或二次型的谱，特别是在“希尔伯特”空间 \(E = \ell_{\mathbf{C}}^{2}(\mathbf{N})\) 上。可以把 E 上任何连续双线性型 B 对应于由 \(\langle Ax, y \rangle = B(x, y)\) 刻画的 E 的连续自同态 A，这一事实为希尔伯特学派所知，尤其是在 E. 施密特、M. 弗雷歇与 F. 里斯 1907–1908 年的工作之后。但正是 F. 里斯在 1913 年一本令人赞叹的书 [59] 中指出了在此背景下把自同态置于首位的益处。

里斯给出了 \(\mathcal{L}(\mathrm{E})\) 的元素 A 的谱的现代定义 [59, p.~139]，指出它是一个有界闭集，并证明预解式 \(\lambda \mapsto (\lambda 1_{\mathrm{E}} - \mathrm{A})^{-1}\) 在 A 的谱的余集上是全纯的。尤其重要的是，他赋予 \(\mathcal{L}(\mathrm{E})\) 的代数结构以中心地位：借鉴沃尔泰拉的工作，他对对称自同态发展了泛函演算的第一个版本，使他能够简单地重新得到希尔伯特与施密特的约化结果。若 A 是 \(\mathcal{L}(\mathrm{E})\) 中的对称元素，里斯说明如何定义 \(f(\mathrm{A})\) ——对每个函数 \(f\) ，它在 \(\mathbf{R}\) 上（下或上）半连续、取值于 \(\mathbf{R}\) ——并指出 \(f \mapsto f(\mathrm{A})\) 就是我们今天所说的代数态射 [59, pp.~129--130]。把这一思想应用于区间 \([\xi, +\infty[\) 的特征函数，就对每个 \(\xi\) 得到 E 的一个对称自同态 \(\mathrm{A}_{\xi}\) ；借助斯蒂尔杰斯积分，里斯随后证明了希尔伯特关系式 (3) 的一个版本：

\[
\langle \mathrm{A} x, y \rangle = \int_ {\mathbf {R}} \xi d \langle \mathrm{A} _ {\xi} x, y \rangle
\]

对 \(x, y \in \mathrm{E}\) 成立，他甚至把这个等式写成

\[
\mathrm{A} = \int_ {\mathbf {R}} \xi d \mathrm{A} _ {\xi}.
\]

于是 A 的谱是这样的值 \(\xi_{0}\) 所成的集合：\(A_{\xi}\) 在 \(\xi_{0}\) 的任何邻域内都不保持常数；而点谱是 \(\xi \mapsto A_{\xi}\) 的不连续点所成的集合。

在以理论的应用作结之前不久，里斯指出 [59, p.~146]，当自同态 A 来自希尔伯特的某个“全连续”双线性型时，其结果具有特别简单的形式：谱约化为特征值所成之集（可能还须加上 0），特征值排成一个趋于 0 的序列，且每个非零特征值的重数是有限的（III, p.~90, 命题 5）。

\histitem{紧算子}

不到三年之后，里斯在其关于泛函方程的论文 [60] 中大大强化了这些评注；该文以历经一个多世纪仍未减损的清晰度，发展了巴拿赫空间上紧算子谱理论的要点。这篇文本宣称的目标是，通过考虑赋有一致收敛拓扑的空间 \(E = \mathcal{C}(I)\) ，对紧区间 I 上的方程 (1) 重新表述弗雷德霍姆理论。但正如我们在早先的一条注记中已经指出的（ÉHM, p.~268），里斯清楚地意识到他的结果适用于任何巴拿赫空间——这一概念要到下一个十年才被引入。

受弗雷歇关于一般拓扑的思想启发，里斯这时考虑 E 的把每个有界序列映为相对紧序列的自同态。与希尔伯特一样，他仍称其为“全连续的”，并指出在 \(\mathrm{E} = \ell_{\mathbf{C}}^{2}(\mathbf{N})\) 的情形，这一新概念等价于他在 1913 年的书中 \(^{(4)}\) 用过的全连续性概念。

里斯研究形如 \(B = 1_{E} - A\) 的自同态，其中 A 是 E 的全连续自同态，并从“局部紧的赋范向量空间必是有限维的”这一——似乎正是为该场合而发现的——事实导出它们的所有谱性质。他容易地证明 B 的核是有限维的，且它的像是闭的、余维数有限。如同 E. 韦尔 [87] 在有限维中所做的那样，他接着考虑迭代 \(B^{k}\) 。他证明它们的核所成的序列与像所成的序列都是平稳的，然后引入我们今天所谓的 B 的零空间与余零空间，并高效地证明 E 是它们的拓扑直和。把这一观察应用于自同态 \(B_{\lambda} = 1_{E} - \lambda A\) （\(\lambda \in C\) ），他便毫不困难地推出巴拿赫空间上全连续自同态的谱性质，而且这几乎是最终的形式。

关于这些算子的谱的主要结果中，似乎只有那些需要对函数空间概念、特别是对偶概念有更成熟认识的结果才超出了里斯的范围。例如，1920 年代末，T. 希尔德布兰特 [33] 与 J. 绍德尔 [63] 将（在此时已充分确立的巴拿赫空间框架内）证明全连续自同态的伴随自同态是全连续的。

此外，里斯的全连续性概念仍然依赖于序列的使用；巴拿赫 [4] 在其 1932 年的书中摆脱了这一点，转而考虑把每个有界集映为相对紧集的映射。“紧”线性映射这一术语要到二十世纪下半叶才逐渐确立，大概是随着 E. 希尔关于算子半群的书 [34] 而流行开来。

1930 年代提出的若干问题长期悬而未决。在这方面，我们注意到 1973 年对“逼近问题”的解决 [18]，巴拿赫与 S. 马祖尔曾在 1932 与 1938 年使这一问题闻名（见 III, p.~16 的注记 6 与 III, p.~112 的习题 25）。

与此同时，洛莫诺索夫 [40] 证明了：在局部凸空间中，与一个非零紧自同态交换的自同态存在非平凡的不变闭子空间（III, p.~13 的推论 2）。这一结果是对任意连续自同态是否存在此类子空间的问题（“不变子空间问题”）最重大的进展之一。该问题在可分希尔伯特空间的连续自同态情形至今仍未解决；P. 恩夫洛 [19] 在 1987 年构造了第一个不具有非平凡不变闭子空间的巴拿赫空间自同态的例子。